\chapter{Geometria afiniczna: równania, ideały i twierdzenie Hilberta o zerach}
\label{ch:afiniczna}
\section{Od układów równań do geometrii}

W geometrii elementarnej opisujemy prostą równaniem $y=mx+b$, parabolę
równaniem $y=x^2$, a okrąg równaniem $x^2+y^2=1$. Geometria algebraiczna
bada analogiczne zbiory w przestrzeniach o dowolnej skończonej liczbie
współrzędnych. Różnica polega na tym, że interesują nas nie tylko
rysunki rozwiązań, ale także wszystkie wielomiany, które na nich znikają,
i funkcje, które można na tych zbiorach określić algebraicznie.

Odpowiemy kolejno na pytania:
\begin{enumerate}[label=\arabic*)]
\item Kiedy dwa układy równań wielomianowych mają ten sam zbiór rozwiązań?
\item Jaką topologię otrzymamy, uznając zbiory rozwiązań za zbiory domknięte?
\item Jak odzyskać informacje o równaniach ze zbioru ich wspólnych zer?
\item Dlaczego odpowiedź zmienia się wraz z ciałem współczynników?
\item Jak przekształcenia geometryczne przekładają się na homomorfizmy pierścieni?
\end{enumerate}

\begin{przyklad}[Te same punkty, różne równania]
W $\C^2$ układy $x=0$ oraz $x^2=0$ mają ten sam zbiór
rozwiązań: oś $y$. Również układy $x=0$ oraz $x^2=0,\ xy=0$
mają ten sam zbiór rozwiązań. Jednak ideały $(x)$, $(x^2)$ oraz
$(x^2,xy)$ nie są jednakowe. Przykładowo w ilorazie
$\C[x,y]/(x^2)$ klasa $x$ jest niezerowa, ale jej kwadrat jest zerowy.
Zatem sam zbiór rozwiązań nie odróżnia wszystkich relacji algebraicznych.
To pierwsza wskazówka, że potrzebujemy jednocześnie zbiorów i ideałów.
\end{przyklad}

\begin{uwaga}[Założenia o ciele]
Przez większość rozdziału $k$ oznacza dowolne \emph{ciało}; jego
charakterystyka może być dodatnia. Przy twierdzeniu Hilberta o zerach
przyjmiemy dodatkowo, że $k$ jest \emph{algebraicznie domknięte}.
Za każdym razem wyraźnie wskażemy, które założenie jest potrzebne.
Wszystkie pierścienie są przemienne z jedynką, a homomorfizmy zachowują jedynkę.
\end{uwaga}

\section{Przestrzeń afiniczna i zbiory algebraiczne}

\begin{definicja}[Przestrzeń afiniczna]
Niech $k$ będzie ciałem, a $n\ge1$ liczbą całkowitą. \emph{Przestrzenią
afiniczną wymiaru $n$ nad $k$}, na poziomie zbioru jej punktów nad $k$,
nazywamy $\A^n(k):=k^n$. Punkt zapisujemy jako
$a=(a_1,\dots,a_n)$, gdzie $a_i\in k$.
\end{definicja}

Symbol $\A^n(k)$ nie oznacza tu jeszcze schematu ani przestrzeni z nadaną
topologią. W tej chwili jest to zbiór uporządkowanych $n$-tek liczb z $k$.
Topologię wprowadzimy dopiero za pomocą równań.

\begin{definicja}[Zbiór zerowy]
Niech $k$ będzie ciałem, $n\ge1$ i $S\subseteq k[x_1,\dots,x_n]$.
Definiujemy \emph{zbiór wspólnych zer wielomianów z $S$} przez
\[
\V_k(S):=\{a\in k^n: f(a)=0\text{ dla każdego }f\in S\}.
\]
Gdy $S=\{f_1,\dots,f_r\}$, piszemy
$\V_k(f_1,\dots,f_r)$, a dla pojedynczego $f$ piszemy $\V_k(f)$.
Podzbiór $X\subseteq k^n$ nazywamy \emph{zbiorem algebraicznym afinicznym},
jeśli istnieje $S\subseteq k[x_1,\dots,x_n]$ takie, że $X=\V_k(S)$.
\end{definicja}

Dopuszczamy $S=\varnothing$. Wtedy każdy punkt spełnia pusty zestaw
równań i $\V_k(\varnothing)=k^n$. Z drugiej strony
$\V_k(1)=\varnothing$, gdyż żadna liczba nie spełnia $1=0$.
Możemy już sprawdzić, dlaczego od listy równań ważniejszy jest ideał.

\begin{stwierdzenie}[Zbiór zerowy zależy tylko od generowanego ideału]
\label{st:V-ideal}
Dla dowolnego ciała $k$, liczby $n\ge1$ i zbioru
$S\subseteq k[x_1,\dots,x_n]$ zachodzi $\V_k(S)=\V_k((S))$,
gdzie $(S)$ jest ideałem generowanym przez $S$.
\end{stwierdzenie}
\begin{proof}
Ponieważ $S\subseteq(S)$, każdy punkt zerujący wszystkie elementy
$(S)$ zeruje też elementy $S$, zatem $\V_k((S))\subseteq\V_k(S)$.
Dla inkluzji odwrotnej ustalmy $a\in\V_k(S)$ i weźmy dowolny
$h\in(S)$. Z definicji ideału generowanego
$h=\sum_{j=1}^r g_js_j$, gdzie $g_j\in k[x_1,\dots,x_n]$,
$s_j\in S$; suma jest skończona. Po podstawieniu punktu $a$ mamy
$h(a)=\sum_jg_j(a)s_j(a)=0$, ponieważ każde $s_j(a)=0$.
Stąd $a\in\V_k((S))$.
\end{proof}

\begin{przyklad}[Zbiór zerowy może być pusty mimo nietrywialnych równań]
W $\R^2$ zbiór $\V_{\R}(x^2+y^2+1)$ jest pusty, bo suma
$x^2+y^2+1$ jest dodatnia w każdym punkcie rzeczywistym.
Tymczasem w $\C^2$ ten sam wielomian ma zera, np. $(i,0)$.
Zmiana ciała współczynników może więc zmienić geometrię zbioru zer.
\end{przyklad}

\begin{przyklad}[Przecięcie dwóch krzywych]
W $\R^2$ układ $y=x^2$ i $y=2-x$ jest równoważny
$y=x^2$, $x^2+x-2=0$. Ponieważ
$x^2+x-2=(x-1)(x+2)$, otrzymujemy dokładnie dwa punkty:
$(1,1)$ i $(-2,4)$. Zapisujemy to jako
$\V_{\R}(y-x^2,y+x-2)=\{(1,1),(-2,4)\}$.
Ten przykład pokazuje, że \emph{przecięcie} zbiorów rozwiązań odpowiada
\emph{połączeniu} ich równań.
\end{przyklad}

\begin{center}
\begin{tikzpicture}[scale=.85,>=Stealth]
\draw[->,gray!75] (-2.9,0)--(2.9,0) node[right] {$x$};
\draw[->,gray!75] (0,-.35)--(0,4.9) node[above] {$y$};
\draw[blue!75!black,line width=1.35pt,domain=-2.15:2.15,samples=100] plot (\x,{\x*\x});
\draw[red!75!black,line width=1.25pt,domain=-2.7:2.35] plot (\x,{2-\x});
\fill (1,1) circle(2pt);
\node[anchor=west,fill=white,inner sep=1pt] at (1.37,1.2) {$(1,1)$};
\fill (-2,4) circle(2pt) node[above right,fill=white,inner sep=1pt] {$(-2,4)$};
\node[blue!75!black] at (1.75,4.65) {$y=x^2$};
\node[red!75!black,fill=white,inner sep=1pt] at (-1.05,3.4) {$y=2-x$};
\end{tikzpicture}
\end{center}
\begin{center}\small Rysunek 1.1. Przecięcie paraboli i prostej nad $\R$; zaznaczono wspólne zera ich równań.\end{center}

\section{Topologia Zariskiego: domknięte są rozwiązania równań}

\subsection{Dlaczego zbiory algebraiczne tworzą rodzinę zbiorów domkniętych?}

Aby zadać topologię na zbiorze $X$, można wskazać rodzinę
$\mathcal C$ jego zbiorów \emph{domkniętych}, pod warunkiem że
$\varnothing,X\in\mathcal C$, dowolne przecięcie zbiorów z
$\mathcal C$ należy do $\mathcal C$, a suma dwóch zbiorów z
$\mathcal C$ także należy do $\mathcal C$. Zbiory otwarte to dopełnienia
tych domkniętych. Sprawdzimy te warunki dla zbiorów zerowych.

\begin{twierdzenie}[Istnienie topologii Zariskiego]
\label{tw:zariski-top}
Niech $k$ będzie dowolnym ciałem i $n\ge1$.
Rodzina $\mathcal C=\{\V_k(S):S\subseteq k[x_1,\dots,x_n]\}$
jest rodziną zbiorów domkniętych pewnej topologii na $k^n$.
Topologię tę nazywamy \emph{topologią Zariskiego}.
Dla ideałów $I,J\subseteq k[x_1,\dots,x_n]$ zachodzą wzory
\[
\V_k(I+J)=\V_k(I)\cap\V_k(J),\qquad
\V_k(IJ)=\V_k(I)\cup\V_k(J).
\]
\end{twierdzenie}
\begin{proof}
Całe $k^n$ jest zbiorem $\V_k(0)$, a zbiór pusty jest równy
$\V_k(1)$.

Niech $\{S_\alpha\}_{\alpha\in T}$ będzie dowolną rodziną zbiorów
wielomianów. Punkt $a$ należy do wszystkich $\V_k(S_\alpha)$
dokładnie wtedy, gdy zeruje wszystkie wielomiany należące do
$\bigcup_{\alpha\in T}S_\alpha$. Dlatego
$\bigcap_{\alpha\in T}\V_k(S_\alpha)=
\V_k(\bigcup_{\alpha\in T}S_\alpha)$.
Z twierdzenia~\ref{st:V-ideal} wynika w szczególności wzór
$\V_k(I+J)=\V_k(I)\cap\V_k(J)$: suma ideałów jest ideałem
generowanym przez zbiór $I\cup J$.

Pozostaje suma dwóch zbiorów. Jeśli $a\in\V_k(I)$, to dla
każdych $f\in I$ i $g\in J$ mamy $f(a)g(a)=0$; każda
skończona suma iloczynów elementów z $I$ i $J$ także zeruje się w $a$.
To daje $\V_k(I)\cup\V_k(J)\subseteq\V_k(IJ)$.

Odwrotnie, niech $a\in\V_k(IJ)$ i załóżmy, że
$a\notin\V_k(I)$. Istnieje wtedy $f\in I$ takie, że
$f(a)\ne0$. Dla dowolnego $g\in J$ wielomian $fg$ należy do
$IJ$, a więc $0=(fg)(a)=f(a)g(a)$. Ponieważ $k$ jest ciałem
i $f(a)\ne0$, dostajemy $g(a)=0$. Dotyczy to każdego $g\in J$,
zatem $a\in\V_k(J)$. Dowiedliśmy odwrotnej inkluzji.
Dzięki domknięciu na dowolne przecięcia i skończone sumy
rodzina $\mathcal C$ jest rodziną zbiorów domkniętych topologii.
\end{proof}

\begin{uwaga}[Dlaczego użyliśmy iloczynu ideałów?]
Równość $\V_k(IJ)=\V_k(I)\cup\V_k(J)$ bierze się z prostego
faktu o liczbach: w ciele $uv=0$ oznacza $u=0$ lub $v=0$.
\emph{Suma} ideałów wymaga jednoczesnego spełnienia dwóch rodzin
równań, więc daje \emph{przecięcie} zbiorów. \emph{Iloczyn} ideałów
wymaga spełnienia co najmniej jednej rodziny, więc daje \emph{sumę}
zbiorów. To odwrócenie intuicji będzie powtarzać się w całym skrypcie.
\end{uwaga}

\begin{przyklad}[Dwie osie i ich punkt wspólny]
W $k[x,y]$ weźmy $I=(x)$ oraz $J=(y)$.
Wtedy $\V_k(I)$ to oś $y$, $\V_k(J)$ to oś $x$,
$\V_k(I+J)=\V_k(x,y)=\{(0,0)\}$, a
$\V_k(IJ)=\V_k(xy)$ jest sumą zbiorów obu osi.
\end{przyklad}

\begin{center}
\begin{tikzpicture}[scale=.82,>=Stealth]
\draw[gray!70,->] (-2.45,0)--(2.45,0) node[right] {$x$};
\draw[gray!70,->] (0,-2.2)--(0,2.4) node[above] {$y$};
\draw[blue!70!black,line width=2pt] (-2.05,0)--(2.05,0);
\draw[red!75!black,line width=2pt] (0,-1.95)--(0,2.0);
\fill[black] (0,0) circle (2.7pt);
\node[blue!70!black,below right,fill=white,inner sep=1pt] at (1.6,-.05) {$\V_k(y)$};
\node[red!75!black,above left,fill=white,inner sep=1pt] at (-.05,1.4) {$\V_k(x)$};
\end{tikzpicture}
\end{center}
\begin{center}\small Rysunek 1.2. Zbiór $\V_{\R}(xy)$: suma zbiorów dwóch przecinających się osi rzeczywistych.\end{center}

\subsection{Najprostsza topologia Zariskiego: prosta afiniczna}

\begin{stwierdzenie}[Zbiory domknięte na prostej nad ciałem nieskończonym]
\label{st:A1-top}
Jeśli $k$ jest nieskończonym ciałem, to domkniętymi zbiorami
Zariskiego w $\A^1(k)$ są dokładnie: zbiór pusty, całe $k$
oraz wszystkie skończone podzbiory $k$.
\end{stwierdzenie}
\begin{proof}
Niech $X=\V_k(S)$ będzie domknięty i różny od całego $k$.
Istnieje wtedy punkt $a\in k\setminus X$. Z definicji $X$
istnieje wielomian $f\in S$ taki, że $f(a)\ne0$;
w szczególności $f$ nie jest wielomianem zerowym.
Każdy punkt $X$ jest pierwiastkiem $f$, zatem $X$ jest skończony:
niezerowy wielomian jednej zmiennej nad ciałem ma nie więcej
pierwiastków niż wynosi jego stopień.

Odwrotnie, dla dowolnego skończonego zbioru
$X=\{a_1,\dots,a_r\}$ mamy
$X=\V_k(\prod_{j=1}^r(x-a_j))$, a dla $X=\varnothing$
można użyć wielomianu stałego $1$. Całe $k$ jest $\V_k(0)$.
\end{proof}

Zbiory otwarte niepuste na tej prostej są więc dopełnieniami
zbiorów skończonych. Dwa dowolne niepuste zbiory otwarte przecinają się,
ponieważ suma dwóch skończonych zbiorów jest skończona, a $k$ jest
nieskończone. W szczególności \emph{topologia Zariskiego na $\A^1(k)$
nie jest Hausdorffa}: nie potrafimy rozdzielić dwóch różnych punktów
rozłącznymi otoczeniami otwartymi.

\begin{przyklad}[Każdy punkt jest domknięty]
Dla dowolnego ciała $k$, liczby całkowitej $n\ge1$
i $a=(a_1,\dots,a_n)\in k^n$
zachodzi $\{a\}=\V_k(x_1-a_1,\dots,x_n-a_n)$.
Topologia Zariskiego spełnia zatem aksjomat $T_1$: wszystkie
jednopunktowe podzbiory są domknięte. Nie przeczy to temu, że
nad ciałem nieskończonym może nie spełniać aksjomatu Hausdorffa.
\end{przyklad}

\begin{przyklad}[Nad ciałem skończonym topologia jest dyskretna]
Niech $n\ge1$. Jeśli $k=\F_q$ ma $q$ elementów, to $k^n$ jest zbiorem skończonym.
Ponieważ każdy punkt jest domknięty, każdy podzbiór $k^n$
stanowi skończoną sumę punktów domkniętych, a zatem też jest domknięty.
W konsekwencji \emph{każdy podzbiór jest również otwarty}.
Topologia Zariskiego na $\A^n(\F_q)$ jest więc dyskretna.
To ostrzeżenie przed przenoszeniem bez zmian intuicji geometrycznej
z $\C$ na ciała skończone.
\end{przyklad}

\subsection{Zbiory podstawowo otwarte}

\begin{definicja}[Zbiór podstawowo otwarty]
Dla ciała $k$, liczby $n\ge1$ oraz $f\in k[x_1,\dots,x_n]$
definiujemy $D_k(f)=\{a\in k^n:f(a)\ne0\}=k^n\setminus\V_k(f)$.
\end{definicja}

\begin{stwierdzenie}[Baza topologii Zariskiego]
Niech $k$ będzie ciałem, $n\ge1$ liczbą całkowitą,
a $f,g\in k[x_1,\dots,x_n]$ wielomianami.
Rodzina $\{D_k(h):h\in k[x_1,\dots,x_n]\}$ jest bazą topologii
Zariskiego na $k^n$. Ponadto $D_k(1)=k^n$, $D_k(0)=\varnothing$
oraz $D_k(f)\cap D_k(g)=D_k(fg)$.
\end{stwierdzenie}
\begin{proof}
Równości $D_k(1)=k^n$ i $D_k(0)=\varnothing$ wynikają z definicji.
Punkt należy do $D_k(f)\cap D_k(g)$ dokładnie wtedy,
gdy $f(a)\ne0$ i $g(a)\ne0$. W ciele jest to równoważne
$(fg)(a)\ne0$, czyli $a\in D_k(fg)$.

Niech $U$ będzie dowolnym zbiorem otwartym.
Jego dopełnienie ma postać $\V_k(S)$ dla pewnego $S$.
Punkt $a$ należy do $U$ wtedy i tylko wtedy, gdy nie spełnia
wszystkich równań z $S$, czyli istnieje $f\in S$ z $f(a)\ne0$.
Zatem $U=\bigcup_{f\in S}D_k(f)$, jak wymaga definicja bazy.
\end{proof}

\begin{uwaga}[Funkcja $1/x$ i otwarta część prostej]
Zbiór $D_k(x)=k\setminus\{0\}$ jest otwarty w $\A^1(k)$.
Na tym zbiorze funkcja $x\mapsto1/x$ ma sens, mimo że nie jest
ograniczeniem wielomianu na całe $k$, gdy $k$ jest nieskończone.
To zapowiada pojęcie \emph{funkcji regularnej na zbiorze otwartym}:
lokalnie będziemy dopuszczać ilorazy wielomianów o nieznikającym
mianowniku. Formalny opis otrzymuje się przez lokalizację pierścienia;
przedstawimy prosty przypadek w dalszej części rozdziału.
\end{uwaga}

\section{Ideał znikania: od geometrii z powrotem do algebry}

\begin{definicja}[Ideał znikania]
Niech $k$ będzie ciałem, $n\ge1$ i $X\subseteq k^n$.
\emph{Ideałem znikania} $X$ nazywamy
\[
\I_k(X)=\{f\in k[x_1,\dots,x_n]:f(a)=0\text{ dla każdego }a\in X\}.
\]
Jeżeli $X=\varnothing$, warunek jest pusty, więc
$\I_k(\varnothing)=k[x_1,\dots,x_n]$.
\end{definicja}

\begin{twierdzenie}[Podstawowe własności ideału znikania]
\label{tw:I-basic}
Niech $k$ będzie ciałem, $n\ge1$ i $X\subseteq k^n$.
Wtedy $\I_k(X)$ jest ideałem radykalnym pierścienia
$k[x_1,\dots,x_n]$. Dla dowolnej rodziny podzbiorów
$\{X_\alpha\}_{\alpha\in T}$ zachodzi
$\I_k(\bigcup_\alpha X_\alpha)=\bigcap_\alpha\I_k(X_\alpha)$.
\end{twierdzenie}
\begin{proof}
Wielomian zerowy znika na $X$, zatem $0\in\I_k(X)$.
Jeżeli $f,g\in\I_k(X)$, to dla dowolnego $a\in X$ mamy
$(f-g)(a)=f(a)-g(a)=0$, więc $f-g\in\I_k(X)$.
Jeśli $h\in k[x_1,\dots,x_n]$, to $(hf)(a)=h(a)f(a)=0$,
a więc $hf\in\I_k(X)$. To dowodzi, że $\I_k(X)$ jest ideałem.

Przypuśćmy, że $f^m\in\I_k(X)$ dla pewnego $m\ge1$.
Dla każdego $a\in X$ zachodzi $f(a)^m=0$ w ciele $k$.
W ciele nie ma niezerowych elementów nilpotentnych, stąd
$f(a)=0$ dla każdego $a\in X$, czyli $f\in\I_k(X)$.
Zatem ideał jest radykalny.

Wreszcie $f$ znika na sumie $\bigcup_\alpha X_\alpha$
dokładnie wtedy, gdy znika na każdym $X_\alpha$ z osobna.
Jest to równoważne $f\in\bigcap_\alpha\I_k(X_\alpha)$.
\end{proof}

\begin{przyklad}[Ideał jednego punktu]
Niech $a=(a_1,\dots,a_n)\in k^n$. W rozdziale o pierścieniach
udowodniliśmy, że jądro homomorfizmu
$\ev_a:k[x_1,\dots,x_n]\to k$, $f\mapsto f(a)$,
wynosi $(x_1-a_1,\dots,x_n-a_n)$.
Ponieważ jądro to dokładnie $\I_k(\{a\})$, mamy
$\I_k(\{a\})=(x_1-a_1,\dots,x_n-a_n)$.
Jest to ideał maksymalny, bo iloraz jest izomorficzny z $k$.
\end{przyklad}

\begin{przyklad}[Rozwiązania układu równań i jego ideał]
Niech $X=\V_k(x,y)\subseteq k^2$. Wtedy $X=\{(0,0)\}$,
a z poprzedniego przykładu $\I_k(X)=(x,y)$.
Natomiast dla $J=(x^2,y^2)$ mamy również
$\V_k(J)=\{(0,0)\}$, lecz $J\ne(x,y)$.
W rzeczywistości $\sqrt J=(x,y)$: elementy $x$ i $y$
należą do $\sqrt J$; zatem $(x,y)\subseteq\sqrt J$.
W drugą stronę, jeśli $f^m\in J$, to $f(0,0)^m=0$,
więc $f(0,0)=0$, a stąd $f\in(x,y)$.
Widać, dlaczego z punktów możemy co najwyżej oczekiwać
odzyskania \emph{pierwiastnika} ideału.
\end{przyklad}

\subsection{Odwracanie inkluzji i operatory domknięcia}

\begin{twierdzenie}[Korespondencja porządku odwracająca inkluzje]
\label{tw:galois}
Niech $k$ będzie ciałem, $n\ge1$, a $R=k[x_1,\dots,x_n]$.
Dla zbiorów $S,T\subseteq R$ i podzbiorów $X,Y\subseteq k^n$
zachodzą implikacje
\begin{align*}
S\subseteq T&\ \Longrightarrow\ \V_k(T)\subseteq\V_k(S),\\
X\subseteq Y&\ \Longrightarrow\ \I_k(Y)\subseteq\I_k(X).
\end{align*}
Ponadto $S\subseteq\I_k(\V_k(S))$ oraz
$X\subseteq\V_k(\I_k(X))$.
W szczególności dla ideału $I\subseteq R$ i zbioru $X\subseteq k^n$
\[
\V_k(\I_k(\V_k(I)))=\V_k(I),\qquad
\I_k(\V_k(\I_k(X)))=\I_k(X).
\]
\end{twierdzenie}
\begin{proof}
Jeżeli $S\subseteq T$, to spełnienie wszystkich równań
z większego zbioru $T$ wymusza spełnienie równań z $S$.
Daje to pierwszą inkluzję. Jeśli $X\subseteq Y$, wielomian
zerujący się na większym zbiorze $Y$ zeruje się również na $X$;
tak otrzymujemy drugą inkluzję.

Każdy wielomian $f\in S$ znika na $\V_k(S)$ z definicji,
czyli $f\in\I_k(\V_k(S))$. Analogicznie każdy punkt $a\in X$
zeruje wszystkie wielomiany, które znikają na $X$, stąd
$a\in\V_k(\I_k(X))$.

Aby dowieść pierwszej z dwóch ostatnich równości, oznaczmy
$Z=\V_k(I)$ i zauważmy, że $I\subseteq\I_k(Z)$, więc
$\V_k(\I_k(Z))\subseteq\V_k(I)=Z$. Inkluzja odwrotna
$Z\subseteq\V_k(\I_k(Z))$ wynika z już udowodnionej własności.
Dowód drugiej równości jest analogiczny: z
$X\subseteq\V_k(\I_k(X))$ wynika
$\I_k(\V_k(\I_k(X)))\subseteq\I_k(X)$,
a odwrotna inkluzja wynika z własności dla $S=\I_k(X)$.
\end{proof}

\begin{center}
\begin{tikzpicture}[>=Stealth,node distance=4cm]
\node[draw,rounded corners,align=center,minimum width=4.7cm,minimum height=1.15cm] (left)
  {rodziny równań\\$I\subseteq k[x_1,\dots,x_n]$};
\node[draw,rounded corners,align=center,right=of left,minimum width=4.7cm,minimum height=1.15cm] (right)
  {zbiory punktów\\$X\subseteq k^n$};
\draw[->,thick,blue!70!black] (left.north east) to[bend left=20]
 node[above] {$I\mapsto\V_k(I)$} (right.north west);
\draw[->,thick,red!70!black] (right.south west) to[bend left=20]
 node[below] {$X\mapsto\I_k(X)$} (left.south east);
\node[below=1.4cm of left,align=center,font=\small] {więcej równań\\$\Longrightarrow$ mniej punktów};
\node[below=1.4cm of right,align=center,font=\small] {więcej punktów\\$\Longrightarrow$ mniej równań};
\end{tikzpicture}
\end{center}
\begin{center}\small Rysunek 1.3. Dwa przejścia odwracające inkluzję: $\V_k$ i $\I_k$.\end{center}

\begin{stwierdzenie}[Domknięcie Zariskiego dowolnego podzbioru]
\label{st:closure}
Dla dowolnego ciała $k$, liczby $n\ge1$ i podzbioru
$X\subseteq k^n$ jego domknięcie w topologii Zariskiego jest równe
$\overline X^{\mathrm{Zar}}=\V_k(\I_k(X))$.
\end{stwierdzenie}
\begin{proof}
Wiemy, że $X\subseteq\V_k(\I_k(X))$, a zbiór po prawej
jest domknięty w topologii Zariskiego.
Niech $Y=\V_k(S)$ będzie dowolnym zbiorem domkniętym zawierającym $X$.
Jeśli $f\in S$, to $f$ znika na $Y$, a więc także na $X$.
Stąd $S\subseteq\I_k(X)$, co odwracając inkluzję daje
$\V_k(\I_k(X))\subseteq\V_k(S)=Y$.
Zatem $\V_k(\I_k(X))$ jest najmniejszym zbiorem domkniętym
zawierającym $X$, czyli jego domknięciem.
\end{proof}

\begin{przyklad}[Nieskończony zbiór jest gęsty na prostej]
Jeżeli $k$ jest nieskończone, a $X\subseteq k$ jest nieskończony,
to $\I_k(X)=(0)$: niezerowy wielomian jednej zmiennej nie może
zerować się na nieskończenie wielu punktach.
Wobec tego $\overline X^{\mathrm{Zar}}=\V_k(0)=k$.
Na przykład $\Z\subseteq\C$ jest gęste w prostej afinicznej
$\A^1(\C)$ w topologii Zariskiego, mimo że w zwykłej topologii
zespolonej $\Z$ jest zbiorem domkniętym i dyskretnym.
\end{przyklad}

\section{Pierścień współrzędnych i algebraiczne funkcje na zbiorze}

\begin{definicja}[Pierścień współrzędnych]
Niech $k$ będzie ciałem, $n\ge1$ i $X\subseteq k^n$ zbiorem
algebraicznym. \emph{Pierścieniem współrzędnych} $X$ nazywamy
$k[X]:=k[x_1,\dots,x_n]/\I_k(X)$.
Klasa $x_i+\I_k(X)$ jest $i$-tą \emph{funkcją współrzędną na $X$}.
\end{definicja}

Konstrukcja jest motywowana następującym rachunkiem. Dwa wielomiany
$f,g$ są równe jako funkcje na $X$ wtedy i tylko wtedy, gdy
$f(a)=g(a)$ dla każdego $a\in X$. To zachodzi dokładnie wtedy,
gdy $(f-g)(a)=0$ dla wszystkich $a\in X$, a więc wtedy i tylko wtedy,
gdy $f-g\in\I_k(X)$. Zatem \emph{elementy $k[X]$ odpowiadają dokładnie
funkcjom wielomianowym na zbiorze punktów $X$}.

\begin{uwaga}[Uwaga o wyjściu od dowolnego ideału]
Jeżeli zaczynamy nie od zbioru $X$, lecz od ideału $J$ i konstruujemy
$k[x_1,\dots,x_n]/J$, to iloraz może zawierać niezerowe nilpotenty.
Natomiast pierścień współrzędnych zbioru punktów $X$ jest zawsze
\emph{zredukowany}, ponieważ $\I_k(X)$ jest radykalny.
Nad ciałem algebraicznie domkniętym okaże się, że
$k[\V_k(J)]\cong k[x_1,\dots,x_n]/\sqrt J$.
Bez tego założenia równość może zawieść.
\end{uwaga}

\begin{przyklad}[Parabola]
Niech $X=\V_k(y-x^2)\subseteq k^2$.
Dla dowolnego ciała $k$ mamy izomorfizm
$k[x,y]/(y-x^2)\cong k[t]$, wynikający z podstawienia
$x=t$, $y=t^2$. Jeśli $k$ jest nieskończone, to
$\I_k(X)=(y-x^2)$: gdy wielomian $f$ znika na wszystkich
punktach $(t,t^2)$, to $f(t,t^2)$ jest wielomianem jednej
zmiennej z nieskończenie wieloma pierwiastkami, czyli jest zerowy;
z poprzedniego rozdziału wiemy, że wówczas $f\in(y-x^2)$.
W takim przypadku $k[X]\cong k[t]$.

Nad $k=\F_q$ mamy natomiast $x^q-x\in\I_k(X)$,
choć $x^q-x\notin(y-x^2)$: gdyby należał, jego obraz w
$k[x,y]/(y-x^2)\cong k[x]$ byłby zerowy, a nie jest.
Nie wolno zatem nad skończonym ciałem utożsamiać obu ilorazów.
\end{przyklad}

\begin{przyklad}[Dwie osie]
Niech $k$ będzie ciałem nieskończonym i
$X=\V_k(xy)\subseteq k^2$. Wtedy $\I_k(X)=(xy)$.
Rzeczywiście, jeśli $f$ znika na osi $y=0$, to
$f(x,0)=0$ jako wielomian w $x$ (bo ma nieskończenie wiele zer),
a zatem $y\mid f$. Analogicznie znikanie na osi $x=0$
daje $x\mid f$. Zapiszmy $f=yg$; ponieważ $f(0,y)=yg(0,y)=0$
jako wielomian, z faktu, że $k[y]$ jest dziedziną, wynika
$g(0,y)=0$, czyli $x\mid g$. Mamy więc $xy\mid f$.
Odwrotna inkluzja $(xy)\subseteq\I_k(X)$ jest oczywista.
Stąd $k[X]\cong k[x,y]/(xy)$. Klasy $x$ oraz $y$
są niezerowe, lecz ich iloczyn jest zerowy: algebraicznie
widzimy dwie składowe przechodzące przez wspólny punkt.
\end{przyklad}

\begin{przyklad}[Krzywa hiperboliczna i odwracalna współrzędna]
Niech $k$ będzie nieskończonym ciałem i
$H=\V_k(xy-1)\subseteq k^2$. Każdy punkt $H$ ma postać
$(t,t^{-1})$ dla pewnego $t\in k^\times:=k\setminus\{0\}$.
W ilorazie $k[x,y]/(xy-1)$ klasy $x$ i $y$ są wzajemnie
odwrotne. Rozważmy \emph{pierścień wielomianów Laurenta}
$k[t,t^{-1}]=\{\sum_{j=-m}^r a_jt^j:a_j\in k\}$.
Homomorfizm $\Phi:k[x,y]\to k[t,t^{-1}]$, zadany przez
$\Phi(x)=t$, $\Phi(y)=t^{-1}$, jest surjektywny:
każda potęga $t^j$ o $j\ge0$ jest obrazem $x^j$, a każda
potęga $t^{-j}$ dla $j>0$ jest obrazem $y^j$.
Ponadto $xy-1\in\kerf\Phi$, więc otrzymujemy homomorfizm
$\overline\Phi:k[x,y]/(xy-1)\to k[t,t^{-1}]$.
Aby wykazać, że jest on izomorfizmem, skonstruujmy odwrotność:
wielomianowi Laurenta $\sum_j a_jt^j$ przyporządkujmy
$\sum_{j\ge0}a_j[x]^j+\sum_{j<0}a_j[y]^{-j}$,
gdzie $[x]$ i $[y]$ oznaczają klasy w ilorazie.
Jest to homomorfizm, ponieważ $[x][y]=1$, zatem
odpowiedniki dodatnich i ujemnych potęg mnożą się zgodnie
z regułą $t^it^j=t^{i+j}$. Złożenia w obu kierunkach
ustalają generatory $t,t^{-1}$ oraz $[x],[y]$, więc są identycznościami.

Obliczmy jeszcze $\I_k(H)$. Jeśli wielomian $f(x,y)$ znika na
wszystkich punktach $(t,t^{-1})$, to wielomian Laurenta
$\Phi(f)$ znika dla wszystkich $t\in k^\times$.
Po pomnożeniu przez dostatecznie wysoką potęgę $t$
otrzymujemy zwykły wielomian jednej zmiennej, który znika
na nieskończonym zbiorze $k^\times$; jest więc wielomianem zerowym.
Stąd $f\in\kerf\Phi=(xy-1)$. Inkluzja odwrotna jest
oczywista, ponieważ $xy-1$ zeruje się na $H$. W rezultacie
$k[H]\cong k[t,t^{-1}]$.
\end{przyklad}

\section{Dlaczego nie każdą nieskończoną rodzinę równań trzeba zapisywać?}

W definicji $\V_k(S)$ dopuściliśmy dowolny, także nieskończony
zbiór wielomianów. Można zapytać, czy taki opis jest rzeczywiście
potrzebny. Odpowiedź daje \emph{twierdzenie Hilberta o bazie}.
Jest to zarazem ważny moment: po raz pierwszy użyjemy własności
pierścienia wielomianów do stwierdzenia czegoś o \emph{wszystkich}
układach równań jednocześnie.

\begin{definicja}[Pierścień noetherowski]
Pierścień przemienny $A$ nazywamy \emph{noetherowskim}, jeśli każdy
jego ideał jest skończenie generowany.
To znaczy: dla każdego ideału $I\subseteq A$ istnieją
$a_1,\dots,a_r\in I$ takie, że $I=(a_1,\dots,a_r)$.
\end{definicja}

\begin{lemat}[Warunek łańcuchów wstępujących]
\label{lem:acc}
Dla pierścienia przemiennego $A$ następujące warunki są równoważne:
\begin{enumerate}[label=\arabic*)]
\item każdy ideał $A$ jest skończenie generowany;
\item każdy wstępujący ciąg ideałów
$I_1\subseteq I_2\subseteq I_3\subseteq\cdots$
jest od pewnego miejsca stały.
\end{enumerate}
\end{lemat}
\begin{proof}
Załóżmy (1) i rozważmy ciąg jak w (2).
Suma mnogościowa $I:=\bigcup_{n\ge1}I_n$ jest ideałem:
jeśli $a\in I_r$ i $b\in I_s$, to dla $m=\max(r,s)$
mamy $a,b\in I_m$, więc $a-b\in I_m\subseteq I$;
analogicznie $ca\in I$ dla $c\in A$.
Jeżeli $I=(0)$, to wszystkie $I_n=(0)$ i ciąg jest stały.
W przeciwnym przypadku z (1) wybieramy niepusty skończony
układ generatorów $a_1,\dots,a_t$ ideału $I$.
Każdy $a_j$ należy do pewnego
$I_{n_j}$. Dla $N=\max_j n_j$ wszystkie generatory należą
do $I_N$, więc $I\subseteq I_N$. Ponieważ $I_N\subseteq I$,
mamy $I=I_N$, a dla $n\ge N$ zachodzi $I_n=I_N$.

Załóżmy teraz (2) i niech $I$ będzie ideałem.
Gdyby $I$ nie był skończenie generowany, moglibyśmy wybrać
$a_1\in I$ i następnie kolejno $a_{r+1}\in I\setminus(a_1,\dots,a_r)$: taki element istnieje, ponieważ żaden
skończony zbiór nie generuje $I$.
Powstałby ściśle rosnący ciąg
$(a_1)\subsetneq(a_1,a_2)\subsetneq\cdots$,
co przeczy (2). Zatem (1) zachodzi.
\end{proof}

\begin{twierdzenie}[Twierdzenie Hilberta o bazie]
\label{tw:hilbert-basis}
Jeśli $A$ jest pierścieniem noetherowskim, to pierścień wielomianów
$A[t]$ również jest noetherowski.
\end{twierdzenie}
\begin{proof}
Niech $I\subseteq A[t]$ będzie dowolnym ideałem.
Dla każdej liczby całkowitej $d\ge0$ zdefiniujmy zbiór
$L_d\subseteq A$ złożony z tych $a\in A$, dla których istnieje
wielomian $f\in I$ o stopniu nie większym niż $d$,
którego współczynnik przy $t^d$ jest równy $a$.
Dopuszczamy więc zerowy współczynnik wiodący (a także wielomian
zerowy). Intuicyjnie $L_d$ zbiera współczynniki, które możemy
\emph{otrzymać na pozycji $t^d$} bez używania wyższych potęg.

Każdy $L_d$ jest ideałem: jeśli $a,b\in L_d$ mają świadków
$f,g\in I$ stopnia co najwyżej $d$, to $f-g\in I$ ma stopień
co najwyżej $d$ i współczynnik $a-b$ przy $t^d$.
Dla $c\in A$ wielomian $cf\in I$ ma taki współczynnik równy $ca$.
Co więcej $L_d\subseteq L_{d+1}$: jeżeli $f$ świadczy, że
$a\in L_d$, to $tf$ ma stopień co najwyżej $d+1$ i współczynnik
$a$ przy $t^{d+1}$. Ponieważ $A$ jest noetherowski,
lemat~\ref{lem:acc} daje liczbę $N\ge0$ taką, że
$L_N=L_{N+1}=L_{N+2}=\cdots$.

Dla każdego $d=0,\dots,N$ wybierzmy skończony układ generatorów
$a_{d,1},\dots,a_{d,r_d}$ ideału $L_d$. Z definicji $L_d$
możemy wybrać wielomiany $f_{d,j}\in I$ stopnia co najwyżej $d$
i ze współczynnikiem $a_{d,j}$ przy $t^d$.
Niech $J$ będzie ideałem $A[t]$ generowanym przez \emph{wszystkie}
$f_{d,j}$, gdzie $0\le d\le N$ i $1\le j\le r_d$.
To skończony zbiór generatorów; oczywiście $J\subseteq I$.

Wykażemy przez indukcję względem stopnia, że każdy $f\in I$
należy do $J$. Dla $f=0$ nie ma czego dowodzić.
Niech $m=\deg f$ i $a$ będzie współczynnikiem przy $t^m$.
Wprowadźmy $d=\min(m,N)$.
Skoro $a\in L_m$, a $L_m=L_N$ gdy $m\ge N$,
to $a\in L_d$. Możemy zatem zapisać
$a=\sum_{j=1}^{r_d}c_ja_{d,j}$ dla pewnych $c_j\in A$.
Wielomian
$h:=\sum_{j=1}^{r_d}c_jt^{m-d}f_{d,j}$ należy do $J$,
ma stopień co najwyżej $m$ i współczynnik $a$ przy $t^m$.
Zatem $f-h\in I$ ma stopień mniejszy od $m$ albo jest zerowy.
Na mocy założenia indukcyjnego $f-h\in J$, a więc $f\in J$.
W rezultacie $I=J$ jest skończenie generowany.
\end{proof}

\begin{wniosek}[Wielomiany w wielu zmiennych]
\label{wn:noetherian-polys}
Jeżeli $k$ jest ciałem i $n\ge1$, to $k[x_1,\dots,x_n]$
jest pierścieniem noetherowskim. W szczególności każdy zbiór
algebraiczny $X\subseteq k^n$ można zapisać w postaci
$X=\V_k(f_1,\dots,f_r)$ dla pewnej skończonej liczby wielomianów.
\end{wniosek}
\begin{proof}
Każdy ideał ciała $k$ jest równy $(0)$ lub $k$:
niezerowy element ideału jest odwracalny, więc taki ideał jest całym
ciałem. Zatem $k$ jest noetherowski.
Stosując twierdzenie~\ref{tw:hilbert-basis} kolejno do
$k[x_1]$, $k[x_1,x_2]=k[x_1][x_2]$ itd., otrzymujemy
noetherowskość $k[x_1,\dots,x_n]$.
Jeśli $X=\V_k(S)$, to ideał $(S)$ ma skończony zbiór
generatorów $f_1,\dots,f_r$. Każdy $f_i$ jest skończoną
sumą wielokrotności wielomianów pochodzących z $S$.
Zbierając wszystkie wielomiany z $S$ występujące w tych
sumach, uzyskujemy skończony podzbiór $S_0\subseteq S$
z własnością $(S_0)=(S)$. Istotnie, $f_i\in(S_0)$ daje
$(S)=(f_1,\dots,f_r)\subseteq(S_0)$, a z $S_0\subseteq S$
wynika inkluzja odwrotna. Zatem
$X=\V_k((S))=\V_k(S_0)$: wystarcza skończenie wiele
\emph{pierwotnych} równań z rodziny $S$.
\end{proof}

\begin{uwaga}[Twierdzenie o istnieniu, nie o algorytmie]
Twierdzenie Hilberta o bazie gwarantuje, że skończony opis istnieje,
ale nie podaje efektywnej metody znalezienia generatorów dla
arbitralnie zadanej rodziny równań. W obliczeniowej geometrii
algebraicznej rolę efektywnego narzędzia pełnią przede wszystkim
\emph{bazy Gröbnera}, które omówimy później tylko wtedy,
gdy okażą się potrzebne do konkretnych rachunków.
\end{uwaga}

\section{Ciała algebraicznie domknięte i problem zer}

\subsection{Dlaczego potrzebujemy domkniętości algebraicznej?}

\begin{definicja}[Ciało algebraicznie domknięte]
Ciało $k$ nazywamy \emph{algebraicznie domkniętym}, jeśli każdy
niezerowy wielomian $f\in k[t]$ dodatniego stopnia ma pierwiastek w $k$.
\end{definicja}

\begin{stwierdzenie}[Rozkład wielomianu na czynniki liniowe]
Niech $k$ będzie ciałem. Następujące warunki są równoważne:
\begin{enumerate}[label=\arabic*)]
\item $k$ jest algebraicznie domknięte;
\item każdy niezerowy wielomian $f\in k[t]$ dodatniego stopnia
rozpada się w $k[t]$ na iloczyn czynników liniowych, z uwzględnieniem
krotności i niezerowego współczynnika stałego;
\item każdy wielomian nierozkładalny w $k[t]$ ma stopień $1$.
\end{enumerate}
\end{stwierdzenie}
\begin{proof}
Załóżmy (1). Dla niezerowego $f$ dodatniego stopnia wybierzmy
pierwiastek $a\in k$. Na mocy twierdzenia o czynniku
$f=(t-a)g$, gdzie $\deg g=\deg f-1$. Powtarzamy tę operację,
a indukcja kończy się na wielomianie stopnia $0$. Otrzymujemy (2).
Jeśli (2), to nierozkładalny wielomian dodatniego stopnia nie może
mieć co najmniej dwóch czynników liniowych, zatem ma stopień $1$.
Tak otrzymujemy (3). Wreszcie każdy wielomian dodatniego stopnia
posiada nierozkładalny dzielnik dodatniego stopnia:
można go znaleźć, rozkładając wielomian, jeśli sam jest rozkładalny,
a następnie postępując indukcyjnie względem stopnia.
Z (3) wynika, że dzielnik ten ma postać $c(t-a)$,
więc $a$ jest pierwiastkiem wyjściowego wielomianu. To daje (1).
\end{proof}

\begin{przyklad}[Przykłady ciał]
Ciało $\C$ jest algebraicznie domknięte (podstawowe twierdzenie algebry).
Ciało $\R$ nie jest: wielomian $t^2+1$ nie ma rzeczywistych pierwiastków.
Ciało $\Q$ też nie jest: wielomian $t^2-2$ nie ma wymiernych
pierwiastków. Ciało skończone $\F_q$ o $q$ elementach nie jest algebraicznie
domknięte, ponieważ wielomian $t^q-t+1$ przyjmuje wartość $1$
na wszystkich jego elementach. Istotnie, dla $a\ne0$
mnożenie przez $a$ permutuje $q-1$ niezerowych elementów ciała.
Mnożąc elementy obu stron tej permutacji, mamy
$\prod_{b\in\F_q^\times}b
=a^{q-1}\prod_{b\in\F_q^\times}b$.
Iloczyn po prawej jest niezerowy, więc $a^{q-1}=1$,
a zatem $a^q=a$. Dla $a=0$ ta równość jest oczywista.
\end{przyklad}

\begin{uwaga}[Algebraiczne domknięcie a nieskończoność]
Każde algebraicznie domknięte ciało jest nieskończone.
Rzeczywiście, gdyby $k=\{a_1,\dots,a_q\}$ było skończone,
wielomian $f(t)=\prod_{i=1}^q(t-a_i)+1$ miałby dodatni stopień,
a dla każdego $a_i$ przyjmowałby wartość $1$.
Nie miałby więc żadnego pierwiastka w $k$.
\end{uwaga}

\subsection{Dwie wersje twierdzenia Hilberta o zerach}

Poniżej udowodnimy dwie fundamentalne wersje
\emph{Nullstellensatz} (twierdzenia Hilberta o zerach).
Pierwsza zapewnia istnienie rozwiązania układu równań;
druga dokładnie identyfikuje wszystkie wielomiany, które
zerują się na zbiorze tych rozwiązań.

\begin{twierdzenie}[Słabe twierdzenie Hilberta o zerach]
\label{tw:weak-NS}
Niech $k$ będzie ciałem algebraicznie domkniętym,
$n\ge1$, a $I\subsetneq k[x_1,\dots,x_n]$ ideałem właściwym.
Wtedy istnieje $a\in k^n$ taki, że $f(a)=0$ dla każdego $f\in I$,
czyli $\V_k(I)\ne\varnothing$.
\end{twierdzenie}

Równoważna kontrapozycja głosi: jeśli układ wielomianowych
równań nad ciałem algebraicznie domkniętym nie ma wspólnego
rozwiązania, to wielomiany generują ideał jednostkowy, czyli
istnieją $f_1,\dots,f_r\in I$ oraz
$h_1,\dots,h_r\in k[x_1,\dots,x_n]$ takie, że
$h_1f_1+\cdots+h_rf_r=1$.
Taką równość można rozumieć jako \emph{algebraiczne świadectwo}
niemożliwości rozwiązania układu.

\begin{twierdzenie}[Silne twierdzenie Hilberta o zerach]
\label{tw:strong-NS}
Niech $k$ będzie ciałem algebraicznie domkniętym,
$n\ge1$ i $I\subseteq k[x_1,\dots,x_n]$ ideałem.
Wtedy
\[
\I_k(\V_k(I))=\sqrt I
=\{f\in k[x_1,\dots,x_n]:f^m\in I\text{ dla pewnego }m\ge1\}.
\]
\end{twierdzenie}

Zauważmy, że tezy dotyczą \emph{każdego} ideału,
a nie tylko ideału głównego czy radykalnego. Najtrudniejszym
składnikiem dowodu słabej wersji będzie lemat Zariskiego.
Udowodnimy go od podstaw, aby nie ukrywać istotnego przejścia
od algebry do geometrii.

\section{Lemat Zariskiego: dlaczego ciało skończenie generowane jako algebra jest algebraiczne?}

\begin{definicja}[Rozszerzenie algebraiczne i element całkowity]
Niech $k\subseteq L$ będzie rozszerzeniem ciał. Element $u\in L$
jest \emph{algebraiczny nad $k$}, jeśli spełnia pewne niezerowe
równanie wielomianowe o współczynnikach w $k$.
Jeżeli $R\subseteq B$ jest zawieraniem pierścieni, to $u\in B$
nazywamy \emph{całkowitym nad $R$}, jeśli istnieją
$a_0,\dots,a_{m-1}\in R$ takie, że
$u^m+a_{m-1}u^{m-1}+\cdots+a_0=0$ dla pewnego $m\ge1$.
W tym drugim przypadku wymagamy, by wielomian był \emph{moniczny},
czyli miał współczynnik wiodący równy $1$.
\end{definicja}

Nad ciałem każda zależność wielomianowa z niezerowym
współczynnikiem wiodącym może zostać podzielona przez ten współczynnik.
Dlatego dla rozszerzenia ciał \emph{algebraiczny} i
\emph{całkowity} oznaczają to samo. Nad ogólnym pierścieniem jest
inaczej: np. $1/2$ jest algebraiczny nad $\Q$, lecz nie jest
całkowity nad $\Z$.

\begin{definicja}[Algebra skończenie generowana]
Niech $k$ będzie ciałem. Pierścień $B$ zawierający $k$
nazywamy \emph{skończenie generowaną $k$-algebrą}, jeśli
istnieją $u_1,\dots,u_s\in B$ takie, że
$B=k[u_1,\dots,u_s]$, czyli każdy element $B$ jest
wartością pewnego wielomianu nad $k$ w punktach $u_1,\dots,u_s$.
Nie należy tego mylić z warunkiem, że $B$ ma skończony wymiar
jako przestrzeń wektorowa nad $k$.
\end{definicja}

\begin{przyklad}[Algebra nie musi mieć skończonego wymiaru]
Pierścień $k[t]$ jest generowany przez jeden element $t$
jako $k$-algebra, ale jego wymiar jako przestrzeni wektorowej
nad $k$ jest nieskończony: wielomiany $1,t,t^2,\dots$
są liniowo niezależne. Lemat Zariskiego mówi, że sytuacja
radykalnie zmienia się, gdy skończenie generowana algebra
jest \emph{ciałem}.
\end{przyklad}

\subsection{Ciało ułamków i algebraiczna niezależność}

Zanim przejdziemy do dowodu lematu Zariskiego, potrzebujemy jeszcze
dwóch pojęć. W pierwszym przechodzimy od dziedziny do ciała,
w którym można dzielić przez każdy element niezerowy.
W drugim sprawdzamy, czy pewne wielkości spełniają jakąkolwiek
relację wielomianową.

\begin{definicja}[Ciało ułamków dziedziny]
Niech $R$ będzie dziedziną całkowitości.
Rozważmy pary $(a,b)$, gdzie $a,b\in R$ i $b\ne0$,
oraz relację $(a,b)\sim(c,d)$ wtedy i tylko wtedy,
gdy $ad=bc$. Klasę pary $(a,b)$ oznaczamy $a/b$.
Zbiór klas nazywamy \emph{ciałem ułamków} $R$
i oznaczamy $\Frac(R)$. Działania definiujemy wzorami
\[
\frac ab+\frac cd=\frac{ad+bc}{bd},\qquad
\frac ab\cdot\frac cd=\frac{ac}{bd}.
\]
\end{definicja}

\begin{lemat}[Poprawność konstrukcji ciała ułamków]
\label{lem:frac-construction}
Dla każdej dziedziny $R$ powyższa relacja jest relacją
równoważności, działania są dobrze określone, a $\Frac(R)$
jest ciałem. Odwzorowanie $R\to\Frac(R)$,
$a\mapsto a/1$, jest różnowartościowym homomorfizmem pierścieni.
\end{lemat}
\begin{proof}
Zwrotność wynika z $ab=ba$, symetria z symetrii równości.
Dla przechodniości załóżmy $a/b=c/d$ i $c/d=e/f$.
Mamy $ad=bc$ i $cf=de$. Mnożąc pierwszą równość przez
$f$, otrzymujemy $adf=bcf=bde$. Ponieważ $d\ne0$
i $R$ jest dziedziną, możemy skrócić $d$; stąd $af=be$,
czyli $a/b=e/f$.

Sprawdźmy niezależność dodawania od reprezentantów.
Jeśli $a/b=a'/b'$ oraz $c/d=c'/d'$, to
$ab'=a'b$ i $cd'=c'd$. Z tych równości otrzymujemy
\begin{align*}
(ad+bc)b'd'
 &=ab'dd'+bb'cd'\\
 &=a'bdd'+bb'c'd\\
 &=(a'd'+b'c')bd.
\end{align*}
Dlatego sumy wyznaczają tę samą klasę. Dla mnożenia
wystarczy pomnożyć równości $ab'=a'b$ oraz $cd'=c'd'$:
$acb'd'=a'c'bd$, czyli również iloczyn nie zależy od
reprezentantów.

Łączność, przemienność i rozdzielność sprowadzają się
po podstawieniu wzorów do odpowiednich własności $R$.
Przykładowo $(a/b)(c/d+e/f)=a(cf+de)/(bdf)$,
podczas gdy $(a/b)(c/d)+(a/b)(e/f)=ac/(bd)+ae/(bf)$;
po sprowadzeniu do wspólnego mianownika $bdf$ otrzymujemy
tę samą klasę $a(cf+de)/(bdf)$.
Elementy neutralne to $0/1$ i $1/1$.
Przeciwnym do $a/b$ jest $(-a)/b$.
Jeśli $a/b\ne0/1$, to $a\ne0$ i $b/a$ jest dozwoloną klasą,
a $(a/b)(b/a)=1/1$. Zatem mamy ciało.
Wreszcie $a/1=0/1$ oznacza $a=0$, więc odwzorowanie
$R\to\Frac(R)$ jest różnowartościowe; zachowanie działań
i jedynki wynika z ich definicji.
\end{proof}

\begin{uwaga}[Ułamki jako funkcje wymierne]
Dla $R=k[t_1,\dots,t_r]$ otrzymujemy ciało
$k(t_1,\dots,t_r)=\Frac(k[t_1,\dots,t_r])$.
Jego elementy mają postać $p(t)/q(t)$, gdzie $p,q$ są
wielomianami, a $q$ jest \emph{niezerowym wielomianem},
niekoniecznie wielomianem, który nie znika w każdym punkcie $k^r$.
To ważne: ułamek jest formalnym elementem ciała ułamków,
nie zaś funkcją określoną we wszystkich punktach przestrzeni afinicznej.
\end{uwaga}

\begin{definicja}[Algebraiczna niezależność]
Niech $k\subseteq L$ będzie rozszerzeniem ciał i
$u_1,\dots,u_r\in L$. Elementy te są
\emph{algebraicznie niezależne nad $k$}, jeśli
$P(u_1,\dots,u_r)\ne0$ dla każdego niezerowego
$P\in k[t_1,\dots,t_r]$. W przeciwnym razie
nazywamy je \emph{algebraicznie zależnymi}.
Pustą rodzinę traktujemy jako algebraicznie niezależną.
\end{definicja}

\begin{przyklad}[Zależność i niezależność]
W ciele $k(t)$ elementy $t$ i $t^2$ są algebraicznie zależne,
ponieważ niezerowy wielomian $P(u,v)=v-u^2$ spełnia
$P(t,t^2)=0$. Natomiast sam element $t$ jest algebraicznie
niezależny nad $k$: gdyby niezerowy wielomian $Q\in k[u]$
spełniał $Q(t)=0$ jako element ciała $k(t)$,
byłby równy zeru jako formalny wielomian, co jest sprzecznością.
\end{przyklad}

\begin{lemat}[Dołączenie elementu do rodziny niezależnej]
\label{lem:independence-algebraic}
Niech $k\subseteq L$ będzie rozszerzeniem ciał,
a $u_1,\dots,u_r\in L$ będą algebraicznie niezależne nad $k$.
Dla $v\in L$ następujące warunki są równoważne:
\begin{enumerate}[label=\arabic*)]
\item rodzina $u_1,\dots,u_r,v$ jest algebraicznie zależna nad $k$;
\item element $v$ jest algebraiczny nad ciałem $k(u_1,\dots,u_r)$.
\end{enumerate}
\end{lemat}
\begin{proof}
Jeśli (1), istnieje niezerowy wielomian
$P(t_1,\dots,t_r,z)=\sum_{j=0}^d P_j(t_1,\dots,t_r)z^j$
z $P(u_1,\dots,u_r,v)=0$.
Przynajmniej jeden $P_j$ jest niezerowym wielomianem.
Algebraiczna niezależność $u_i$ oznacza, że
$P_j(u_1,\dots,u_r)\ne0$, więc otrzymujemy
niezerowy wielomian jednej zmiennej $z$ o współczynnikach
w ciele $k(u_1,\dots,u_r)$, który zeruje się dla $z=v$.
To znaczy, że zachodzi (2).

Odwrotnie, jeśli $v$ jest algebraiczny nad
$F=k(u_1,\dots,u_r)$, to
$\sum_{j=0}^d c_j v^j=0$ dla pewnych $c_j\in F$,
nie wszystkich zerowych. Wszystkie $c_j$ są ilorazami
wielomianów w $u_1,\dots,u_r$. Po pomnożeniu przez wspólny
niezerowy mianownik otrzymujemy niezerową zależność
$\sum_{j=0}^d Q_j(u_1,\dots,u_r)v^j=0$
z $Q_j\in k[t_1,\dots,t_r]$, nie wszystkimi zerowymi.
Zatem rodzina $u_1,\dots,u_r,v$ jest algebraicznie zależna.
\end{proof}

\begin{lemat}[Skończona liczba elementów całkowitych]
\label{lem:integral-finite}
Niech $R\subseteq B$ będzie zawieraniem pierścieni i niech
elementy $u_1,\dots,u_s\in B$ będą całkowite nad $R$.
Wtedy $R[u_1,\dots,u_s]$ jest skończenie generowanym
$R$-modułem: istnieją elementy $b_1,\dots,b_m$
tej algebry takie, że każdy jej element można zapisać jako
$\sum_{j=1}^m r_jb_j$ z $r_j\in R$.
\end{lemat}
\begin{proof}
Najpierw rozważmy jeden element $u$, który spełnia
$u^d+a_{d-1}u^{d-1}+\cdots+a_0=0$.
Przekształcamy tę równość do postaci
$u^d=-a_{d-1}u^{d-1}-\cdots-a_0$.
Każdą wyższą potęgę $u^N$ możemy kolejno zredukować,
mnożąc tę relację przez $u^{N-d}$. Indukcja względem $N$
daje $R[u]=R\cdot1+R\cdot u+\cdots+R\cdot u^{d-1}$.

Dla kilku elementów $u_j$ wybierzmy dla każdego moniczne
równanie stopnia $d_j$. Każdy iloczyn potęg tych elementów,
a zatem dowolny element $R[u_1,\dots,u_s]$,
można zredukować do kombinacji monomów
$u_1^{e_1}\cdots u_s^{e_s}$ z $0\le e_j<d_j$.
Jest tylko $d_1\cdots d_s$ takich monomów.
\end{proof}

\begin{lemat}[Dlaczego elementy skończonej algebry spełniają równania moniczne]
\label{lem:determinant-integral}
Niech $R\subseteq B$ będą pierścieniami i załóżmy, że $B$
jest skończenie generowanym $R$-modułem. Wtedy każdy element
$z\in B$ jest całkowity nad $R$.
\end{lemat}
\begin{proof}
Wybierzmy generatory $b_1,\dots,b_m$ modułu $B$ nad $R$.
Mnożenie przez $z$ przenosi każdy generator w element $B$,
możemy więc znaleźć współczynniki $a_{ij}\in R$ takie, że
$zb_i=\sum_{j=1}^m a_{ij}b_j$.
Niech $M=(a_{ij})\in M_m(R)$ i niech
$b=(b_1,\dots,b_m)^T$. Otrzymujemy równanie
$(zE-M)b=0$, gdzie $E$ jest macierzą jednostkową.

Dla macierzy kwadratowej $C$ przez $\operatorname{adj}(C)$
oznaczamy transponowaną macierz dopełnień algebraicznych.
Rozwinięcie wyznacznika względem wiersza lub kolumny daje
$\operatorname{adj}(C)C=\det(C)E$: na przekątnej otrzymujemy
wyznacznik $C$, a poza przekątną wyznacznik macierzy
z dwoma jednakowymi wierszami, czyli zero.
Tożsamość obowiązuje też nad dowolnym pierścieniem przemiennym.
Stosujemy ją do $C=zE-M$ i dostajemy
$\operatorname{adj}(zE-M)(zE-M)=\det(zE-M)E$.
Po pomnożeniu równania przez tę macierz dostajemy
$\det(zE-M)b_i=0$ dla każdego $i$.
Ponieważ generatory $b_i$ generują cały $B$ jako $R$-moduł,
a w szczególności generują element $1\in B$,
wynika stąd $\det(zE-M)\cdot1=0$ w $B$.
Wielomian $P(T):=\det(TE-M)\in R[T]$ jest moniczny
stopnia $m$. Otrzymaliśmy $P(z)=0$, czyli $z$ jest całkowity nad $R$.
\end{proof}

\begin{lemat}[Dziedzina całkowitości nad pierścieniem]
\label{lem:field-integral}
Niech $R\subseteq L$ będą pierścieniami, gdzie $L$ jest ciałem.
Jeżeli każdy element $L$ jest całkowity nad $R$, to $R$ jest ciałem.
\end{lemat}
\begin{proof}
Ustalmy $0\ne r\in R$. W ciele $L$ istnieje $r^{-1}$.
Na mocy założenia spełnia ono równanie
$(r^{-1})^m+a_{m-1}(r^{-1})^{m-1}+\cdots+a_0=0$
z $a_j\in R$. Mnożąc przez $r^{m-1}$, otrzymujemy
$ r^{-1}=-(a_{m-1}+a_{m-2}r+\cdots+a_0r^{m-1})\in R$.
Zatem każdy niezerowy element $R$ ma odwrotność w $R$.
\end{proof}

\begin{lemat}[Lokalizacja pierścienia wielomianów nie staje się ciałem]
\label{lem:localization-not-field}
Niech $k$ będzie ciałem, $r\ge1$ i
$0\ne d\in k[t_1,\dots,t_r]$.
Pierścień $k[t_1,\dots,t_r,1/d]$, złożony z wyrażeń
$p/d^m$ z $p\in k[t_1,\dots,t_r]$ i $m\ge0$,
nie jest ciałem.
\end{lemat}
\begin{proof}
Wprowadźmy $P=k[t_1,\dots,t_r]$ oraz $s=1+t_1d\in P$.
Wielomian $s$ jest niezerowy i niestały: stopień iloczynu
$t_1d$ względem $t_1$ jest o jeden większy od stopnia $d$,
więc dodanie stałej $1$ nie może usunąć wyrazu najwyższego stopnia.
Wielomian niestały nie jest jednostką pierścienia $P$,
ponieważ stopnie niezerowych iloczynów w zmiennej $t_1$
sumują się, a dla zmiennych pozostałych współczynniki tworzą
dziedzinę $k[t_2,\dots,t_r]$.

Załóżmy dla sprzeczności, że $s$ jest odwracalny w
$P[1/d]$. Wtedy istnieją $p\in P$ i $m\ge0$ takie, że
$s(p/d^m)=1$, czyli $sp=d^m$ w dziedzinie $P$.
W ilorazie $P/(s)$ zachodzi $t_1d=-1$.
Oznacza to, że klasa $d$ ma odwrotność $-t_1$.
Jednocześnie z $sp=d^m$ wynika, że $d^m=0$ w tym ilorazie.
Jeśli $m\ge1$, odwracalność $d$ daje $1=0$ w $P/(s)$.
Jeśli $m=0$, z $sp=1$ otrzymujemy to samo.
W obu przypadkach $1\in(s)$, więc $s$ jest jednostką $P$,
sprzecznie z poprzednim akapitem.
\end{proof}

\begin{twierdzenie}[Lemat Zariskiego]
\label{tw:zariski-lemma}
Niech $k$ będzie ciałem, a $L$ ciałem zawierającym $k$.
Jeśli istnieją $a_1,\dots,a_n\in L$ takie, że
$L=k[a_1,\dots,a_n]$ jako $k$-algebra, to każdy element $L$
jest algebraiczny nad $k$. W rzeczywistości $L$ jest
skończenie wymiarową przestrzenią wektorową nad $k$.
\end{twierdzenie}
\begin{proof}
Spośród generatorów $a_1,\dots,a_n$ wybierzmy maksymalną
(ze względu na liczbę elementów) podrodzinę algebraicznie
niezależną nad $k$. Po zmianie numeracji możemy przyjąć,
że jest nią $a_1,\dots,a_r$; być może $r=0$.
Algebraiczna niezależność oznacza, że dla każdego
niezerowego $P\in k[t_1,\dots,t_r]$ mamy
$P(a_1,\dots,a_r)\ne0$.
Dlatego $k[a_1,\dots,a_r]$ jest izomorficzny z
pierścieniem wielomianów $P=k[t_1,\dots,t_r]$,
a wewnątrz ciała $L$ możemy rozważyć jego ciało ułamków
$F=k(a_1,\dots,a_r)$.

Z maksymalności wybranej rodziny i
lematu~\ref{lem:independence-algebraic} wynika, że każdy
pozostały generator $a_j$, $j>r$, jest algebraiczny nad $F$.
Gdyby któryś nie był, moglibyśmy go dołączyć do rodziny,
zachowując algebraiczną niezależność, co przeczyłoby jej
maksymalności. Wobec tego dla każdego $j>r$ istnieje
równanie moniczne
$ a_j^{m_j}+c_{j,m_j-1}a_j^{m_j-1}+\cdots+c_{j,0}=0$
o współczynnikach $c_{j,l}\in F$.
Każdy taki współczynnik jest ilorazem dwóch elementów $P$.
Wybierzmy jeden niezerowy $d\in P$ podzielny przez wszystkie
mianowniki występujące w skończonej liczbie tych równań.
Wtedy $c_{j,l}\in R:=P[1/d]$, więc wszystkie $a_j$ dla $j>r$
są całkowite nad $R$.

Zauważmy, że $R\subseteq L$: elementy $P$ są w $L$,
a $1/d$ również należy do $L$, bo $L$ jest ciałem i $d\ne0$.
Co więcej,
$L=k[a_1,\dots,a_n]=R[a_{r+1},\dots,a_n]$:
lewa strona zawiera prawą, zaś prawa zawiera wszystkie generatory
lewej. Na mocy lematu~\ref{lem:integral-finite}
$L$ jest skończenie generowanym $R$-modułem.
Z lematu~\ref{lem:determinant-integral} każdy element $L$
jest całkowity nad $R$, a z lematu~\ref{lem:field-integral}
wynika, że $R$ musi być ciałem.

Gdyby $r\ge1$, przeczyłoby to lematowi
\ref{lem:localization-not-field}, ponieważ
$R=k[t_1,\dots,t_r,1/d]$. Zatem $r=0$.
Oznacza to, że wszystkie generatory $a_j$ są algebraiczne nad
samym $k$. Możemy przyjąć $R=k$ i zastosować ponownie
lemat~\ref{lem:integral-finite}: otrzymujemy, że
$L=k[a_1,\dots,a_n]$ jest skończenie wymiarowe nad $k$.
Każdy element takiego rozszerzenia jest algebraiczny
(np. na mocy lematu~\ref{lem:determinant-integral}), co kończy dowód.
\end{proof}

\begin{uwaga}[Gdzie tkwi idea dowodu?]
Gdyby $L$ zawierało element transcendentalny $t$, to konieczne
byłyby odwrotności wielu wielomianów w $t$ (na przykład
$1/(1+t)$), ponieważ $L$ jest ciałem.
Tymczasem skończona liczba generatorów algebry pozwala
odwrócić tylko mianowniki pochodzące z pewnego skończonego
zbioru wielomianów. Element $1+t d$ z dowodu pokazuje,
że zawsze znajdziemy nowy nieodwracalny mianownik.
To jest geometryczno-algebraiczny mechanizm stojący za lematem Zariskiego.
\end{uwaga}

\section{Dowód twierdzenia Hilberta o zerach}

\subsection{Od maksymalnego ideału do punktu}

\begin{lemat}[Istnienie ideału maksymalnego w pierścieniu noetherowskim]
\label{lem:maximal-exists}
Niech $R$ będzie niezerowym pierścieniem noetherowskim.
Wtedy istnieje ideał maksymalny $\mathfrak m\subsetneq R$.
Co więcej, każdy ideał właściwy $I\subsetneq R$ jest zawarty
w pewnym ideale maksymalnym.
\end{lemat}
\begin{proof}
Rozważmy rodzinę wszystkich ideałów właściwych zawierających $I$.
Nie jest ona pusta, ponieważ zawiera sam $I$.
Gdyby żaden element tej rodziny nie był maksymalny ze względu na
inkluzję, moglibyśmy rozpocząć od $I_0=I$ i wybierać kolejno
ideały właściwe $I_0\subsetneq I_1\subsetneq I_2\subsetneq\cdots$.
Jest to sprzeczne z warunkiem łańcuchów wstępujących
z lematu~\ref{lem:acc}. Zatem istnieje ideał właściwy
$\mathfrak m$ zawierający $I$ i nieposiadający większego
ideału właściwego. Z definicji jest maksymalny.
Dla pierwszej części wystarczy przyjąć $I=(0)$;
jest to ideał właściwy, ponieważ $R\ne(0)$.
\end{proof}

\begin{lemat}[Ciało skończenie generowane nad ciałem algebraicznie domkniętym]
\label{lem:extension-trivial}
Niech $k$ będzie algebraicznie domkniętym ciałem,
a $L$ ciałem zawierającym $k$.
Jeśli $L$ jest skończenie generowaną $k$-algebrą,
to $L=k$.
\end{lemat}
\begin{proof}
Na mocy lematu Zariskiego każdy $a\in L$ jest algebraiczny
nad $k$. Wybierzmy niezerowy wielomian
$f\in k[t]$ o najmniejszym stopniu spośród wielomianów
spełniających $f(a)=0$. Ponieważ $k$ jest algebraicznie
domknięte, $f$ rozpada się na czynniki liniowe:
$f(t)=c\prod_{j=1}^r(t-b_j)$, gdzie $b_j\in k$ i $c\ne0$.
Z równości
$0=f(a)=c\prod_{j=1}^r(a-b_j)$ w ciele $L$
wynika, że $a=b_j$ dla pewnego $j$, a więc $a\in k$.
Dowolny element $L$ należy do $k$, stąd $L=k$.
\end{proof}

\begin{proof}[Dowód słabego twierdzenia Hilberta o zerach~\ref{tw:weak-NS}]
Oznaczmy $R=k[x_1,\dots,x_n]$ i niech $I\subsetneq R$.
Na mocy wniosku~\ref{wn:noetherian-polys} pierścień $R$
jest noetherowski; pierścień ilorazowy $R/I$ też jest
noetherowski, ponieważ jego ideały odpowiadają ideałom $R$
zawierającym $I$, a ilorazy generatorów tych ideałów generują
odpowiednie ideały w $R/I$.
Ponieważ $I$ jest właściwy, $R/I$ jest niezerowy.
Z lematu~\ref{lem:maximal-exists} istnieje ideał maksymalny
$\overline{\mathfrak m}$ w $R/I$.
Jego przeciwobraz względem rzutowania $R\to R/I$
jest ideałem maksymalnym $\mathfrak m$ w $R$ zawierającym $I$:
rzeczywiście $R/\mathfrak m\cong(R/I)/\overline{\mathfrak m}$,
a ostatni iloraz jest ciałem.

Rozważmy $L:=R/\mathfrak m$.
Jest to ciało oraz skończenie generowana $k$-algebra:
$L=k[\overline x_1,\dots,\overline x_n]$,
gdzie $\overline x_i=x_i+\mathfrak m$.
Włączenie $k\to L$ jest różnowartościowe: jego jądro jest
ideałem ciała $k$, a ponieważ $1\mapsto1\ne0$,
jądro jest zerowe. Możemy więc traktować $k$ jako podciało $L$.
Na mocy lematu~\ref{lem:extension-trivial} mamy $L=k$.
Każda klasa $\overline x_i$ odpowiada zatem pewnemu
$a_i\in k$. Niech $a=(a_1,\dots,a_n)\in k^n$.

Dla każdego $f\in R$ jego klasa w $L$ jest równa
$f(\overline x_1,\dots,\overline x_n)$, a po identyfikacji
$L=k$ jest to $f(a_1,\dots,a_n)$.
Jeżeli $f\in I\subseteq\mathfrak m$, jego klasa w $L$
jest zerowa, więc $f(a)=0$.
Punkt $a$ jest wspólnym zerem wszystkich wielomianów z $I$.
\end{proof}

\begin{wniosek}[Maksymalne ideały są ideałami punktów]
\label{wn:maximal-point}
Jeżeli $k$ jest algebraicznie domkniętym ciałem i $n\ge1$,
to każdy ideał maksymalny $\mathfrak m\subseteq k[x_1,\dots,x_n]$
ma dokładnie jedną postać
$\mathfrak m=(x_1-a_1,\dots,x_n-a_n)$ z $a\in k^n$.
\end{wniosek}
\begin{proof}
Zastosujmy dowód słabego twierdzenia do samego ideału
maksymalnego $\mathfrak m$.
Otrzymamy punkt $a$ taki, że wszystkie wielomiany z
$\mathfrak m$ znikają w $a$, czyli
$\mathfrak m\subseteq\I_k(\{a\})=(x_1-a_1,\dots,x_n-a_n)$.
Ideał po prawej jest właściwy, bo wielomian $1$ nie znika w $a$.
Maksymalność $\mathfrak m$ wymusza równość.
Dla jednoznaczności: jeśli ideały punktów $a$ i $b$ są równe,
to wielomian $x_i-a_i$ należy do ideału punktu $b$,
więc $0=(x_i-a_i)(b)=b_i-a_i$ dla każdego $i$.
Stąd $a=b$.
\end{proof}

\begin{uwaga}[Gdzie użyliśmy domkniętości algebraicznej?]
Jedyny moment w dowodzie, gdy potrzebowaliśmy algebraicznej
domkniętości $k$, to przejście od skończenie generowanego
\emph{ciała} $L=R/\mathfrak m$ do równości $L=k$.
Nad $\R$ możemy mieć $L\cong\C$, jak w przykładzie
$\R[t]/(t^2+1)$. Wówczas ideał maksymalny nie musi być
ideałem rzeczywistego punktu.
\end{uwaga}

\subsection{Trik Rabinowitscha i silne twierdzenie Hilberta o zerach}

\begin{lemat}[Trik Rabinowitscha]
\label{lem:rabinowitsch}
Niech $k$ będzie ciałem algebraicznie domkniętym, $n\ge1$,
$I\subseteq k[x_1,\dots,x_n]$ ideałem i
$0\ne f\in k[x_1,\dots,x_n]$ wielomianem, który znika
we wszystkich punktach $\V_k(I)$.
Wtedy istnieje $N\ge1$ takie, że $f^N\in I$.
\end{lemat}
\begin{proof}
Niech $R=k[x_1,\dots,x_n]$.
Z twierdzenia Hilberta o bazie możemy wybrać generatory
$I=(g_1,\dots,g_r)$. Dodajmy nową zmienną $t$ i
rozważmy ideał
$J=(g_1,\dots,g_r,1-tf)\subseteq R[t]$.
Załóżmy, że punkt $(a,b)\in k^{n+1}$ zeruje wszystkie
generatory $J$. Równości $g_j(a)=0$ oznaczają,
że $a\in\V_k(I)$, a więc z założenia $f(a)=0$.
Ostatnie równanie $1-bf(a)=0$ staje się jednak $1=0$,
co jest niemożliwe. Zatem $\V_k(J)=\varnothing$.

Słabe twierdzenie Hilberta o zerach, zastosowane w $n+1$
zmiennych, daje $J=R[t]$.
Zatem istnieją $u_1,\dots,u_r,v\in R[t]$ takie, że
\[
1=\sum_{j=1}^r u_j(x,t)g_j(x)+v(x,t)(1-tf(x)).
\]
Ponieważ $R$ jest dziedziną, możemy przejść do ciała jego
ułamków $K=\Frac(R)$, w którym niezerowy wielomian $f$
ma odwrotność. Podstawmy $t=1/f$ do tożsamości w $K$.
Ostatni składnik znika, ponieważ $1-(1/f)f=0$, i otrzymujemy
$1=\sum_{j=1}^r u_j(x,1/f)g_j(x)$.

Każde $u_j$ jest wielomianem w $t$ o współczynnikach z $R$.
Po podstawieniu $t=1/f$ jego współczynniki są ilorazami
wielomianów przez potęgi $f$. Wybieramy $N\ge1$
nie mniejsze od wszystkich występujących potęg mianownika.
Po pomnożeniu równości przez $f^N$ otrzymujemy
$f^N=\sum_{j=1}^r h_j(x)g_j(x)$ z $h_j\in R$.
A więc $f^N\in I$, jak należało wykazać.
\end{proof}

\begin{proof}[Dowód silnego twierdzenia Hilberta o zerach~\ref{tw:strong-NS}]
Oznaczmy $R=k[x_1,\dots,x_n]$.
Najpierw wykażmy $\sqrt I\subseteq\I_k(\V_k(I))$.
Jeśli $f\in\sqrt I$, to istnieje $m\ge1$ takie,
że $f^m\in I$. Dla dowolnego $a\in\V_k(I)$ mamy
$f(a)^m=f^m(a)=0$. Ponieważ $k$ jest ciałem,
$f(a)=0$, czyli $f\in\I_k(\V_k(I))$.
Ta inkluzja zachodzi nawet bez założenia algebraicznej domkniętości.

Odwrotnie, niech $f\in\I_k(\V_k(I))$.
Jeżeli $f=0$, to $f\in\sqrt I$ bez dalszych obliczeń.
Jeśli $f\ne0$, to lemat~\ref{lem:rabinowitsch}
daje $N\ge1$ z $f^N\in I$, czyli $f\in\sqrt I$.
To kończy dowód równości.
\end{proof}

\begin{przyklad}[Algebraiczne świadectwo braku punktów]
Rozważmy w $\C[x,y]$ wielomiany $f_1=x$ oraz
$f_2=xy-1$. Wspólne zera nie istnieją: z $x=0$
wynika $xy-1=-1\ne0$.
Równość $1=yx-(xy-1)$ pokazuje bezpośrednio,
że $(x,xy-1)=(1)$.
Słabe twierdzenie Hilberta o zerach gwarantuje,
że dla \emph{dowolnego} układu wielomianów bez wspólnych
zer nad $\C$ istnieje podobne świadectwo wielomianowe.
\end{przyklad}

\begin{przyklad}[Dlaczego powstaje potęga?]
Niech $I=(x^3,x^2y)\subseteq\C[x,y]$.
Wówczas $\V_{\C}(I)=\V_{\C}(x)$, bo $x^3=0$
w ciele daje $x=0$. Wielomian $x$ znika na tym zbiorze,
ale $x\notin I$: każda niezerowa suma wielokrotności
generatorów ma wszystkie monomy podzielne przez $x^2$.
Jednak $x^3\in I$, więc $x\in\sqrt I$.
Silne twierdzenie mówi dokładnie, że zawsze wystarczy
pewna dodatnia potęga znikającego wielomianu.
\end{przyklad}

\section{Co dokładnie daje Nullstellensatz?}

\begin{twierdzenie}[Korespondencja ideałów radykalnych i zbiorów algebraicznych]
\label{tw:radical-correspondence}
Niech $k$ będzie ciałem algebraicznie domkniętym, $n\ge1$
i $R=k[x_1,\dots,x_n]$.
Odwzorowania $I\mapsto\V_k(I)$ i $X\mapsto\I_k(X)$
ustanawiają wzajemnie odwrotne bijekcje pomiędzy rodziną
ideałów radykalnych $I\subseteq R$ a rodziną
zbiorów algebraicznych $X\subseteq k^n$.
Bijekcje odwracają inkluzje.
\end{twierdzenie}
\begin{proof}
Dla ideału radykalnego $I$ silne twierdzenie Hilberta o zerach
daje $\I_k(\V_k(I))=\sqrt I=I$.
Jeśli $X$ jest zbiorem algebraicznym, to $X=\V_k(J)$
dla pewnego ideału $J$. Z twierdzenia~\ref{tw:galois}
wiemy, że $\V_k(\I_k(\V_k(J)))=\V_k(J)$,
a więc $\V_k(\I_k(X))=X$.
Każdy $\I_k(X)$ jest radykalny na mocy
tego, co wykazaliśmy w twierdzeniu~\ref{tw:I-basic}.
Otrzymujemy obie potrzebne identyczności.
Odwracanie inkluzji było dowiedzione w twierdzeniu~\ref{tw:galois}.
\end{proof}

\begin{center}
\begin{tikzpicture}[>=Stealth]
\node[draw,rounded corners,fill=gray!5,align=center,minimum width=4.1cm,minimum height=1cm] (a) at (0,0)
 {ideały radykalne\\$I=\sqrt I$};
\node[draw,rounded corners,fill=gray!5,align=center,minimum width=4.1cm,minimum height=1cm] (b) at (7,0)
 {zbiory algebraiczne\\$X=\V_k(\I_k(X))$};
\draw[->,blue!70!black,thick] (a.north east) to[bend left=19]
 node[above] {$\V_k$} (b.north west);
\draw[->,red!75!black,thick] (b.south west) to[bend left=19]
 node[below] {$\I_k$} (a.south east);
\node[font=\small] at (3.5,-1.35)
 {$k$ algebraicznie domknięte};
\end{tikzpicture}
\end{center}
\begin{center}\small Rysunek 1.4. Korespondencja ideałów i zbiorów punktów po przejściu do ideałów radykalnych.\end{center}

\begin{wniosek}[Pierścień współrzędnych odczytany z równań]
\label{wn:coordinate-ring}
Niech $k$ będzie ciałem algebraicznie domkniętym,
$n\ge1$ i $I\subseteq k[x_1,\dots,x_n]$ ideałem.
Dla $X=\V_k(I)$ istnieje naturalny izomorfizm
$k[X]\cong k[x_1,\dots,x_n]/\sqrt I$.
Jeśli $I$ jest radykalny, to $k[X]\cong k[x_1,\dots,x_n]/I$.
\end{wniosek}
\begin{proof}
Z definicji $k[X]=k[x_1,\dots,x_n]/\I_k(X)$.
Z silnego twierdzenia Hilberta o zerach mamy
$\I_k(X)=\I_k(\V_k(I))=\sqrt I$. Podstawiamy tę równość.
\end{proof}

\begin{przyklad}[Nierówność obiektów nad $\R$]
Niech $I=(x^2+1)\subseteq\R[x]$.
Iloraz $\R[x]/I$ jest ciałem izomorficznym z $\C$, więc
$I$ jest ideałem maksymalnym, a zatem radykalnym.
Mimo to $\V_{\R}(I)=\varnothing$, skąd
$\I_{\R}(\V_{\R}(I))=\R[x]$, nie zaś $I$.
Jest to kontrprzykład do silnego Nullstellensatz nad
ciałem, które nie jest algebraicznie domknięte.
\end{przyklad}

\begin{przyklad}[Nierówność obiektów nad $\F_q$]
Niech $k=\F_q$, a $I=(0)\subseteq k[x]$.
Pierścień $k[x]$ jest dziedziną, więc $I$ jest radykalny.
Jednak $\V_k(I)=k$ i niezerowy wielomian $x^q-x$
znika we wszystkich punktach $k$. Stąd
$x^q-x\in\I_k(\V_k(I))$, mimo że $x^q-x\notin\sqrt I=(0)$.
Samo założenie, że $k$ jest ciałem, zdecydowanie nie wystarcza.
\end{przyklad}

\section{Nierozkładalność zbioru a pierwszość ideału}

\begin{definicja}[Zbiór nierozkładalny]
Niech $T$ będzie przestrzenią topologiczną, a $X\subseteq T$
niepustym zbiorem domkniętym.
Mówimy, że $X$ jest \emph{nierozkładalny}, jeśli dla każdych
dwóch podzbiorów domkniętych $Y,Z\subseteq X$ z
$X=Y\cup Z$ zachodzi $Y=X$ lub $Z=X$.
Jeśli istnieje rozkład na dwa właściwe zbiory domknięte,
to $X$ jest \emph{rozkładalny}.
\end{definicja}

Nierozkładalność nie oznacza tego samego co spójność topologiczna.
Zbiór nierozkładalny jest spójny, ponieważ rozkład na dwa
rozłączne niepuste zbiory jednocześnie otwarte i domknięte
dawałby w szczególności rozkład na dwa właściwe zbiory domknięte.
Odwrotna implikacja jest na ogół fałszywa.

\begin{twierdzenie}[Nierozkładalność a ideały pierwsze]
\label{tw:irreducible-prime}
Niech $k$ będzie dowolnym ciałem, $n\ge1$
i $X\subseteq k^n$ niepustym zbiorem algebraicznym.
Wtedy $X$ jest nierozkładalny w topologii Zariskiego
wtedy i tylko wtedy, gdy $\I_k(X)$ jest ideałem pierwszym.
Równoważnie: $X$ jest nierozkładalny wtedy i tylko wtedy,
gdy jego pierścień współrzędnych $k[X]$ jest dziedziną całkowitości.
\end{twierdzenie}
\begin{proof}
Najpierw załóżmy, że $X$ jest nierozkładalny.
Niech $f,g\in k[x_1,\dots,x_n]$ oraz $fg\in\I_k(X)$.
Każdy punkt $X$ spełnia $f(a)g(a)=0$,
a więc należy do $\V_k(f)$ albo do $\V_k(g)$.
Otrzymujemy rozkład
$X=(X\cap\V_k(f))\cup(X\cap\V_k(g))$
na dwa podzbiory domknięte w $X$. Nierozkładalność oznacza,
że jeden z nich jest całym $X$.
Jeśli $X\subseteq\V_k(f)$, to $f\in\I_k(X)$;
w przeciwnym razie $g\in\I_k(X)$. Ponieważ $X\ne\varnothing$,
ideał $\I_k(X)$ jest właściwy: wielomian $1$ nigdzie nie znika.
Wykazaliśmy więc, że $\I_k(X)$ jest pierwszy.
Ta implikacja nie wymaga algebraicznej domkniętości $k$.

Odwrotnie, załóżmy, że $\mathfrak p:=\I_k(X)$ jest pierwszy,
a $X=Y\cup Z$ z $Y,Z$ domkniętymi w $X$.
Ponieważ $X$ jest domknięty w $k^n$, również $Y$ i $Z$
są domknięte w $k^n$. Możemy zatem zapisać
$Y=\V_k(J)$ i $Z=\V_k(K)$ dla pewnych ideałów $J,K$.
Wtedy $X=\V_k(JK)$.
Każdy element $JK$ znika na $X$, więc $JK\subseteq\I_k(X)=\mathfrak p$.
Jeśli $J\not\subseteq\mathfrak p$, wybierzmy $f\in J\setminus\mathfrak p$.
Dla dowolnego $g\in K$ iloczyn $fg\in JK\subseteq\mathfrak p$,
zatem pierwszość $\mathfrak p$ wymusza $g\in\mathfrak p$.
Stąd $K\subseteq\mathfrak p$ i
$Z=\V_k(K)\supseteq\V_k(\mathfrak p)$.
Ponieważ $X$ jest zbiorem algebraicznym, mamy
$\V_k(\mathfrak p)=\V_k(\I_k(X))=X$
na mocy twierdzenia~\ref{tw:galois}, bez dodatkowych
założeń o ciele $k$. Zatem $Z=X$.
Jeśli zaś $J\subseteq\mathfrak p$, otrzymujemy
$Y\supseteq\V_k(\mathfrak p)=X$, więc $Y=X$.
W obu przypadkach rozkład jest trywialny.

Ostatnia równoważność wynika z kryterium ilorazowego:
ideał $\I_k(X)$ jest pierwszy wtedy i tylko wtedy,
gdy $k[x_1,\dots,x_n]/\I_k(X)=k[X]$ jest dziedziną.
\end{proof}

\begin{uwaga}[Dlaczego rozróżniamy dwa kryteria?]
Dla każdego ciała $k$ mamy równoważność:
$X$ niepusty algebraiczny nierozkładalny wtedy i tylko wtedy,
gdy \emph{jego ideał znikania} $\I_k(X)$ jest pierwszy.
Inne, silniejsze zdanie brzmi: dla dowolnego ideału pierwszego
$\mathfrak p\subseteq k[x_1,\dots,x_n]$ zbiór
$\V_k(\mathfrak p)$ jest niepusty i nierozkładalny.
To drugie zdanie wymaga dodatkowych założeń. Nad $\R$
ideał pierwszy $(x^2+1)$ ma pusty zbiór zer; nad ciałem
skończonym $\F_q$ ideał $(0)\subseteq\F_q[x]$
jest pierwszy, ale jego zbiór zer $\F_q$ ma dyskretną
topologię i jest rozkładalny, jeśli $q>1$.
\end{uwaga}

\begin{przyklad}[Parabola nierozkładalna nad $\C$]
Niech $X=\V_{\C}(y-x^2)$.
Wtedy $\I_{\C}(X)=(y-x^2)$, a iloraz
$\C[x,y]/(y-x^2)\cong\C[t]$ jest dziedziną.
Wobec tego $X$ jest nierozkładalny.
Jest to ścisłe uzasadnienie geometrycznej interpretacji
izomorfizmu, który pojawił się w rozdziale o pierścieniach.
\end{przyklad}

\begin{przyklad}[Suma dwóch osi jest rozkładalna]
Niech $X=\V_{\C}(xy)$.
Mamy $X=\V_{\C}(x)\cup\V_{\C}(y)$,
a oba zbiory po prawej są właściwymi podzbiorami domkniętymi
$X$. Ponadto $\C[X]\cong\C[x,y]/(xy)$ ma dzielniki zera:
klasy $x$ i $y$ są niezerowe, a ich iloczyn jest zerem.
Widzimy ten sam fenomen z perspektywy topologii i algebry.
\end{przyklad}

\begin{twierdzenie}[Skończony rozkład na składowe nierozkładalne]
\label{tw:finite-components}
Niech $k$ będzie dowolnym ciałem, $n\ge1$ i
$X\subseteq k^n$ niepustym zbiorem algebraicznym.
Wtedy istnieją niepuste nierozkładalne podzbiory
algebraiczne $X_1,\dots,X_r\subseteq X$ takie, że
$X=X_1\cup\cdots\cup X_r$.
Można wybrać je jako \emph{maksymalne nierozkładalne podzbiory}
$X$; po usunięciu powtórzeń wybór ten jest jednoznaczny.
\end{twierdzenie}
\begin{proof}
Najpierw pokażemy, że w topologii Zariskiego nie istnieje
nieskończony ściśle malejący ciąg zbiorów algebraicznych
$Y_1\supsetneq Y_2\supsetneq\cdots$.
Gdyby taki istniał, ideały znikania tworzyłyby ciąg
$\I_k(Y_1)\subseteq\I_k(Y_2)\subseteq\cdots$.
Inkluzje byłyby ściśle: jeśli $\I_k(Y_i)=\I_k(Y_{i+1})$,
to po zastosowaniu $\V_k$ dostalibyśmy
$Y_i=\V_k(\I_k(Y_i))=\V_k(\I_k(Y_{i+1}))=Y_{i+1}$,
sprzeczność. Ale pierścień $k[x_1,\dots,x_n]$ jest noetherowski,
więc taki nieskończenie rosnący ciąg ideałów nie istnieje.

Załóżmy teraz, że istnieje niepusty zbiór algebraiczny,
który \emph{nie} daje się zapisać jako skończona suma zbiorów
nierozkładalnych. Spośród wszystkich takich zbiorów wybierzmy
minimalny względem inkluzji: wybór jest możliwy dzięki
właśnie udowodnionemu warunkowi braku nieskończonych ściśle
malejących ciągów. Niech ten zbiór nazywa się $Y$.
Z definicji kontrprzykładu $Y$ nie jest nierozkładalny,
a więc $Y=Y'\cup Y''$ dla właściwych domkniętych
podzbiorów $Y',Y''\subsetneq Y$.
Z minimalności $Y$ oba zbiory $Y',Y''$ mają skończone
rozkłady nierozkładalne. Łącząc te rozkłady, otrzymujemy
taki rozkład również dla $Y$, sprzeczność.

Mamy więc skończony rozkład $X=Z_1\cup\cdots\cup Z_s$.
Jeśli któryś $Z_i$ zawiera się w innym $Z_j$, możemy go usunąć.
Powtarzając tę czynność, otrzymujemy rozkład bez zawierań.
Każdy pozostały $Z_i$ jest maksymalnym nierozkładalnym
podzbiorem $X$: gdyby istniał nierozkładalny
$Z_i\subsetneq W\subseteq X$, to z
$W=\bigcup_j(W\cap Z_j)$ i nierozkładalności $W$
wynikałoby $W\subseteq Z_j$ dla pewnego $j$.
Zatem $Z_i\subsetneq W\subseteq Z_j$, co przeczy brakowi zawierań.

Aby uzasadnić jednoznaczność, zauważmy najpierw, że jeśli
nierozkładalny zbiór $W$ zawiera się w skończonej sumie
domkniętych zbiorów $Z_1\cup\cdots\cup Z_s$,
to $W\subseteq Z_j$ dla pewnego $j$: dla $s=2$ jest to definicja,
a dla większego $s$ wynika przez indukcję po $s$.
Niech $X=W_1\cup\cdots\cup W_t$ będzie drugim rozkładem
na maksymalne nierozkładalne zbiory. Każdy $Z_i$ zawiera się
w pewnym $W_j$ na mocy poprzedniego faktu, a maksymalność
$Z_i$ wymusza $Z_i=W_j$. Zamieniając role rozkładów,
otrzymujemy tę samą rodzinę podzbiorów z dokładnością
do kolejności.
\end{proof}

\begin{przyklad}[Trzy proste]
Zbiór $X=\V_{\C}(xy(x-y))\subseteq\C^2$
jest sumą zbiorów trzech prostych: $x=0$, $y=0$ oraz $x=y$.
Każda jest nierozkładalna, bo jej pierścień współrzędnych
jest izomorficzny z $\C[t]$. Żadna nie zawiera innej,
więc są to dokładnie trzy składowe nierozkładalne $X$.
\end{przyklad}

\section{Morfizmy afiniczne i odwrócenie kierunku odwzorowań}

\subsection{Wielomianowe odwzorowania zbiorów algebraicznych}

\begin{definicja}[Odwzorowanie regularne zbiorów afinicznych]
Niech $k$ będzie ciałem algebraicznie domkniętym, $n,m\ge1$,
a $X\subseteq k^n$ i $Y\subseteq k^m$ niepustymi zbiorami algebraicznymi.
Odwzorowanie $F:X\to Y$ nazywamy \emph{regularnym} (w sensie
afinicznym), jeśli istnieją wielomiany
$f_1,\dots,f_m\in k[x_1,\dots,x_n]$ takie, że
$F(a)=(f_1(a),\dots,f_m(a))$ dla każdego $a\in X$.
Wielomiany $f_i$ nie muszą być jednoznaczne, ale ich klasy w $k[X]$
są jednoznacznie określone przez $F$.
\end{definicja}

\begin{przyklad}[Parametryzacja paraboli]
Niech $k$ będzie ciałem algebraicznie domkniętym i $X=\A^1(k)=k$ i
$Y=\V_k(y-x^2)\subseteq k^2$.
Odwzorowanie $F:k\to Y$, $t\mapsto(t,t^2)$, jest regularne,
bo jego współrzędne są wielomianami. Odwzorowanie odwrotne
$G:Y\to k$, $(x,y)\mapsto x$, jest również regularne,
a złożenia $G\circ F$ i $F\circ G$ są identycznościami.
Parabola jest więc \emph{izomorficzna jako zbiór algebraiczny}
z prostą afiniczną. Sam kształt paraboli na rysunku nie jest
tutaj niezmiennikiem geometrycznym: ważna jest struktura funkcji
wielomianowych.
\end{przyklad}

\begin{definicja}[Homomorfizm $k$-algebr]
Niech $A,B$ będą pierścieniami zawierającymi ciało $k$ jako podpierścień.
Homomorfizm pierścieni $\varphi:A\to B$ nazywamy
\emph{homomorfizmem $k$-algebr}, gdy $\varphi(c)=c$
dla każdego $c\in k$.
\end{definicja}

\begin{twierdzenie}[Kontrawariantny opis odwzorowań regularnych]
\label{tw:affine-morphism}
Niech $k$ będzie ciałem algebraicznie domkniętym, $n,m\ge1$,
a $X\subseteq k^n$ i $Y\subseteq k^m$ niepustymi zbiorami
algebraicznymi. Wtedy każdemu odwzorowaniu regularnemu
$F:X\to Y$ odpowiada homomorfizm $k$-algebr
\[
F^*:k[Y]\longrightarrow k[X],\qquad
F^*([h])=[h\circ F],
\]
gdzie $[h]$ oznacza klasę funkcji wielomianowej.
To przyporządkowanie jest bijekcją między
odwzorowaniami regularnymi $X\to Y$ a homomorfizmami
$k$-algebr $k[Y]\to k[X]$.
Jeżeli ponadto $q\ge1$, $Z\subseteq k^q$ jest niepustym
zbiorem algebraicznym, a $G:Y\to Z$ odwzorowaniem regularnym,
to $(G\circ F)^*=F^*\circ G^*$.
\end{twierdzenie}
\begin{proof}
Niech $F=(f_1,\dots,f_m)$.
Jeśli $[h]=[h']$ w $k[Y]$, to wielomiany $h,h'$ mają
te same wartości w każdym punkcie $Y$. Ponieważ $F(a)\in Y$
dla każdego $a\in X$, otrzymujemy
$h(F(a))=h'(F(a))$ dla wszystkich $a\in X$.
Zatem $[h\circ F]=[h'\circ F]$ w $k[X]$ i wzór na $F^*$
nie zależy od reprezentanta.
Podstawienie wielomianów zachowuje sumę, iloczyn i stałe,
więc $F^*$ jest homomorfizmem $k$-algebr.

W przeciwnym kierunku niech
$\varphi:k[Y]\to k[X]$ będzie homomorfizmem $k$-algebr.
Oznaczmy przez $[y_1],\dots,[y_m]$ klasy funkcji współrzędnych
na $Y$. Dla każdego $j$ wybierzmy wielomian
$f_j\in k[x_1,\dots,x_n]$ reprezentujący
$\varphi([y_j])\in k[X]$.
Zdefiniujmy $F:X\to k^m$ przez
$F(a)=(f_1(a),\dots,f_m(a))$.
Musimy sprawdzić, że faktycznie $F(X)\subseteq Y$.
Dla dowolnego $h\in\I_k(Y)$ klasa $[h]$ w $k[Y]$ jest zerowa.
Stąd $\varphi([h])=0$ w $k[X]$.
Ponieważ $h$ jest wielomianem w zmiennych $y_1,\dots,y_m$,
mamy $\varphi([h])=[h(f_1,\dots,f_m)]$.
Ta klasa jest zerowa dokładnie wtedy, gdy
$h(f_1(a),\dots,f_m(a))=0$ dla wszystkich $a\in X$.
Oznacza to, że każdy punkt $F(a)$ zeruje wszystkie wielomiany
ideału $\I_k(Y)$. Ponieważ $Y$ jest algebraiczny,
$Y=\V_k(\I_k(Y))$, a więc $F(a)\in Y$.
Otrzymaliśmy odwzorowanie regularne $F:X\to Y$.

Skonstruowany homomorfizm $F^*$ jest równy $\varphi$,
ponieważ oba mają te same wartości na stałych z $k$
i na klasach $[y_j]$, które generują $k[Y]$ jako $k$-algebrę.
Z drugiej strony odwzorowanie regularne $F$ jest jednoznacznie
wyznaczone przez jego $m$ funkcji współrzędnych na $X$.
Obie konstrukcje są więc odwrotne.

Dla złożenia i dowolnego $h\in k[Z]$ mamy
$((G\circ F)^*h)(a)=h(G(F(a)))$
oraz $(F^*(G^*h))(a)=(G^*h)(F(a))=h(G(F(a)))$.
Zatem $(G\circ F)^*=F^*\circ G^*$.
\end{proof}

\begin{center}
\begin{tikzpicture}[>=Stealth]
\node (x) at (0,1.3) {$X$};
\node (y) at (5.5,1.3) {$Y$};
\draw[->,thick,blue!70!black] (x)--node[above] {$F$}(y);
\node (kx) at (0,-.55) {$k[X]$};
\node (ky) at (5.5,-.55) {$k[Y]$};
\draw[->,thick,red!70!black] (ky)--node[below] {$F^*$}(kx);
\draw[densely dashed,gray!70] (0,.95)--(0,-.16);
\draw[densely dashed,gray!70] (5.5,.95)--(5.5,-.16);
\end{tikzpicture}
\end{center}
\begin{center}\small Rysunek 1.5. Odwzorowanie punktów $F:X\to Y$ indukuje homomorfizm pierścieni w kierunku przeciwnym.\end{center}

\begin{uwaga}[Dlaczego strzałka biegnie w drugą stronę?]
Jeśli mamy funkcję $h$ określoną na $Y$ i odwzorowanie
$F:X\to Y$, możemy złożyć je jako $h\circ F$,
otrzymując funkcję na $X$. Innymi słowy, funkcje na \emph{celu}
można cofać do funkcji na \emph{źródle}.
Wiele fundamentalnych konstrukcji geometrii algebraicznej
ma właśnie taki kontrawariantny charakter.
\end{uwaga}

\subsection{Zaskakujący przykład: bijekcja nie musi być izomorfizmem}

\begin{przyklad}[Krzywa z ostrzem i jej parametryzacja]
\label{pr:cusp}
Niech $k$ będzie ciałem algebraicznie domkniętym.
Rozważmy krzywą
$C=\V_k(y^2-x^3)\subseteq k^2$ oraz odwzorowanie
$\nu:k\to C$, $t\mapsto(t^2,t^3)$.
Jest ono regularne i rzeczywiście przyjmuje wartości w $C$,
bo $(t^3)^2-(t^2)^3=0$.

Pokażemy, że $\nu$ jest bijekcją na zbiorach punktów.
Jeżeli $(x,y)\in C$ i $x=0$, to $y^2=0$, a więc $y=0$;
wówczas jedynym parametrem jest $t=0$.
Jeśli $x\ne0$, definiujemy $t=y/x$. Z równania $y^2=x^3$
wynika $t^2=y^2/x^2=x$, a następnie
$t^3=t\,t^2=(y/x)x=y$. Zatem każdy punkt krzywej ma postać
$(t^2,t^3)$. Jednoznaczność także wynika z rachunku:
dla punktu różnego od $(0,0)$ parametr musi wynosić $y/x$.

Mimo to $\nu$ \emph{nie jest izomorfizmem algebraicznym}.
Aby tego dowieść, obliczymy homomorfizm pierścieni
$\Phi:k[x,y]\to k[t]$, $\Phi(x)=t^2$, $\Phi(y)=t^3$.
Ma on obraz $k[t^2,t^3]$.
Pokażemy, że jego jądro jest równe $(y^2-x^3)$.
Dowolny wielomian $f(x,y)$ możemy zredukować modulo
$y^2-x^3$ do postaci $A(x)+yB(x)$:
za każdym razem zastępujemy $y^m$, $m\ge2$, przez
$x^3y^{m-2}$, używając równania $y^2=x^3$.
Jeśli $\Phi(f)=0$, to po podstawieniu mamy
$A(t^2)+t^3B(t^2)=0$.
Pierwsza suma zawiera tylko parzyste potęgi $t$,
a druga tylko nieparzyste potęgi $t$ stopnia co najmniej $3$.
Nie mogą się więc wzajemnie znosić. Wszystkie współczynniki
obu wielomianów są zerowe, czyli $A=B=0$.
To znaczy $f\in(y^2-x^3)$.
Ponieważ ciało algebraicznie domknięte $k$ jest nieskończone,
dla dowolnego wielomianu $f$ warunek, że $f$ znika na
całym $C$, jest (dzięki surjektywności $\nu$) równoważny temu,
że wielomian $\Phi(f)(t)=f(t^2,t^3)$ znika dla każdego
$t\in k$, czyli jest zerowy jako wielomian.
Stąd $\I_k(C)=\kerf\Phi=(y^2-x^3)$.
W konsekwencji pierwsze twierdzenie o izomorfizmie daje
\[
k[C]\cong k[x,y]/(y^2-x^3)\cong k[t^2,t^3]\subsetneq k[t].
\]
Ostatnia inkluzja jest ścisła, ponieważ $t\notin k[t^2,t^3]$:
każdy monom o dodatnim stopniu w generatorach $t^2,t^3$
ma stopień co najmniej $2$, zatem z ich skończonej sumy
nie otrzymamy wielomianu stopnia $1$.
Gdyby $\nu$ miała regularną odwrotność, odpowiadający jej
homomorfizm $\nu^*:k[C]\to k[t]$ musiałby być izomorfizmem,
a nie jest nawet surjektywny.
\end{przyklad}

\begin{center}
\begin{tikzpicture}[>=Stealth,scale=.78]
\draw[->,gray!75] (-.25,0)--(4.8,0) node[right] {$x$};
\draw[->,gray!75] (0,-3.3)--(0,3.4) node[above] {$y$};
\draw[blue!75!black,very thick,domain=-1.63:1.63,samples=120,smooth]
 plot ({\x*\x},{\x*\x*\x});
\fill (0,0) circle (2.3pt);
\draw[densely dashed] (1,0)--(1,1);
\fill[red!75!black] (1,1) circle(2.3pt);
\node[right,red!75!black] at (1,1) {$(t^2,t^3)$ dla $t=1$};
\node[blue!75!black] at (3.1,2.8) {$y^2=x^3$};
% Objaśnienie przeniesiono do podpisu, aby nie przysłaniało dolnej gałęzi.
\end{tikzpicture}
\end{center}
\begin{center}\small Rysunek 1.6. Rzeczywisty obraz krzywej z ostrzem; nad $\C$ krzywa ma również punkty zespolone.\end{center}

\begin{uwaga}[Czym różni się ten przykład od paraboli?]
Dla paraboli parametr $t$ jest samą współrzędną $x$,
więc można go odzyskać wielomianowo. Na krzywej z ostrzem
parametr poza początkiem jest równy $y/x$.
Jest to funkcja wymierna, której zapis ma w początku
zerowy mianownik. Nie istnieje globalny opis tego parametru
przez element pierścienia $k[C]$.
W rozdziale o osobliwościach wrócimy do tego przykładu:
$\nu$ jest prototypem \emph{normalizacji} krzywej.
\end{uwaga}

\subsection{Obraz odwzorowania regularnego nie musi być domknięty}

\begin{przyklad}[Rzut hiperboli na oś]
Niech $k$ będzie algebraicznie domkniętym ciałem i
$H=\V_k(xy-1)\subseteq k^2$.
Rozważmy rzut $p:H\to k$, $p(x,y)=x$.
Jest to odwzorowanie regularne. Jego obraz jest równy
$k^\times=k\setminus\{0\}$: punkt hiperboli ma $x\ne0$,
a dla każdego $x\ne0$ punkt $(x,x^{-1})$ leży na $H$.

Zbiór $k\setminus\{0\}$ nie jest domknięty w topologii
Zariskiego na $k$, ponieważ $k$ jest nieskończone,
a jedyne właściwe zbiory domknięte na prostej są skończone.
Obraz odwzorowania regularnego może więc być otwarty, ale niedomknięty.
To ważne ostrzeżenie przed bezkrytycznym przenoszeniem własności
odwzorowań właściwych z topologii na wszystkie morfizmy algebraiczne.
\end{przyklad}

\begin{center}
\begin{tikzpicture}[>=Stealth,scale=.79]
\draw[->,gray!75] (-3.8,0)--(3.9,0) node[right] {$x$};
\draw[->,gray!75] (0,-3.25)--(0,3.4) node[above] {$y$};
\draw[blue!70!black,very thick,domain=.32:3.1,samples=100]
 plot (\x,{1/\x});
\draw[blue!70!black,very thick,domain=-3.1:-.32,samples=100]
 plot (\x,{1/\x});
\draw[->,red!70!black,thick] (2,0.5)--(2,0.08);
\draw[->,red!70!black,thick] (-2,-0.5)--(-2,-0.08);
\draw[fill=white,thick] (0,0) circle(4pt);
\node[blue!70!black] at (2.75,2.2) {$xy=1$};
% Objaśnienie rzutu jest w podpisie rysunku.
\end{tikzpicture}
\end{center}
\begin{center}\small Rysunek 1.7. Rzut punktów hiperboli na oś współrzędnych pomija zero.\end{center}

\section{Kolejne intuicje: wymiar i kierunek dalszej teorii}

\begin{definicja}[Wymiar przez łańcuchy nierozkładalnych podzbiorów]
Niech $k$ będzie ciałem algebraicznie domkniętym,
$n\ge1$, a $X\subseteq k^n$ zbiorem algebraicznym.
Dla $X\ne\varnothing$ definiujemy \emph{wymiar} $X$ jako kres górny liczb
całkowitych $r\ge0$, dla których istnieją
nierozkładalne, niepuste, domknięte podzbiory
$X_0\subsetneq X_1\subsetneq\cdots\subsetneq X_r\subseteq X$.
Zapisujemy $\dim X$. Gdy $X=\varnothing$, przyjmujemy
umownie $\dim X=-\infty$.
\end{definicja}

\begin{przyklad}[Punkt ma wymiar zero]
Jeśli $X=\{a\}$ jest pojedynczym punktem, to jedynym
niepustym podzbiorem domkniętym $X$ jest $X$.
Zatem istnieje łańcuch długości $0$ (samo $X$),
ale nie istnieje łańcuch długości $1$.
Stąd $\dim\{a\}=0$.
\end{przyklad}

\begin{przyklad}[Prosta afiniczna ma wymiar jeden]
Niech $k$ będzie algebraicznie domkniętym ciałem.
Ponieważ $k$ jest nieskończone, właściwe domknięte
podzbiory $\A^1(k)$ są skończone.
Cała prosta jest nierozkładalna: nie można zapisać
nieskończonego $k$ jako sumy dwóch skończonych zbiorów.
Wybierzmy dowolne $a\in k$. Łańcuch
$\{a\}\subsetneq\A^1(k)$ pokazuje, że wymiar jest co najmniej $1$.
Nie istnieje jednak dłuższy łańcuch nierozkładalnych
podzbiorów: każdy właściwy niepusty nierozkładalny zbiór
domknięty na prostej jest singletonem, ponieważ skończony
zbiór o co najmniej dwóch punktach można rozłożyć na dwa
właściwe domknięte podzbiory.
Zatem $\dim\A^1(k)=1$.
\end{przyklad}

Dla $\A^n(k)$ zachodzi $\dim\A^n(k)=n$, ale pełne
udowodnienie tego faktu wymaga dalszych narzędzi algebry przemiennej
(np. normalizacji Noethera lub teorii wysokości ideałów).
Nie będziemy zastępować brakującego dowodu intuicją z rysunku.
Szczegółową teorię wymiaru wprowadzimy przed badaniem
osobliwości i krzywych algebraicznych.

\section{Podsumowanie: sieć zależności między pojęciami}

Oto logika rozdziału. Dla dowolnego ciała $k$ i
$R=k[x_1,\dots,x_n]$ możemy przejść od zbioru równań
$S\subseteq R$ do zbioru ich zer $\V_k(S)$, a od zbioru
punktów $X\subseteq k^n$ do ideału znikania $\I_k(X)$.
Każde z tych przejść odwraca inkluzje. Ponadto
$\overline X^{\mathrm{Zar}}=\V_k(\I_k(X))$,
a $\I_k(X)$ jest zawsze radykalny.

Gdy \emph{dodatkowo} $k$ jest algebraicznie domknięte,
silne twierdzenie Hilberta o zerach gwarantuje
$\I_k(\V_k(I))=\sqrt I$ dla każdego ideału $I$.
Dlatego zbiory algebraiczne odpowiadają dokładnie
ideałom radykalnym, a nie wszystkim ideałom.
Z kolei zbiory nierozkładalne odpowiadają ideałom pierwszym.
Pierścienie współrzędnych przenoszą geometrię na grunt algebry,
a odwzorowaniom regularnym odpowiadają homomorfizmy
pierścieni biegnące w przeciwnym kierunku.
% ======================================================================
% TRZECI ROZDZIAŁ CAŁOŚCI: NUMER 2, PONIEWAŻ SKRYPT ZACZYNA SIĘ OD 0.
% ======================================================================
